\documentclass[12pt,oneside]{book}

% =====================================================
% METADATOS DEL DOCUMENTO
% =====================================================
\newcommand{\universidad}{Universidad de Buenos Aires (UBA)}
\newcommand{\facultad}{Facultad de Derecho}
\newcommand{\carrera}{Cátedra de Lectocomprensión --- Abogacía / Traductorado}
\newcommand{\materia}{Lectocomprensión de Textos Jurídicos en Inglés}
\newcommand{\titulo}{Manual integrado de estudio}
\newcommand{\autor}{Isaac Edgar Camacho Ocampo}
\newcommand{\anio}{2026}

% =====================================================
% PAQUETES
% =====================================================
\usepackage[utf8]{inputenc}
\usepackage[T1]{fontenc}
\usepackage[provide=*,spanish]{babel}

\usepackage[
    top=2.5cm,
    bottom=2.5cm,
    left=3cm,
    right=2.5cm,
    headheight=15pt
]{geometry}

\usepackage{amsmath}
\usepackage{amssymb}
\usepackage{mathtools}

\usepackage{graphicx}
\usepackage{float}
\usepackage{caption}
\usepackage{subcaption}

\usepackage{booktabs}
\usepackage{array}
\usepackage{longtable}
\usepackage{tabularx}

\usepackage{enumitem}

\usepackage{xcolor}
\usepackage[most]{tcolorbox}

\usepackage{fancyhdr}
\usepackage{ragged2e}

\usepackage[
    colorlinks=true,
    linkcolor=blue!50!black,
    citecolor=green!40!black,
    urlcolor=blue!60!black,
    pdftitle={\materia{} --- \titulo},
    pdfauthor={\autor},
    pdfsubject={Apuntes universitarios integrados},
    bookmarks=true
]{hyperref}

% =====================================================
% RUTAS DE IMÁGENES
% =====================================================
\graphicspath{
    {./img/}
    {./graficos/}
    {./figuras/}
}

% =====================================================
% CONFIGURACIÓN GENERAL
% =====================================================
\setcounter{tocdepth}{2}

\setlength{\parindent}{0pt}
\setlength{\parskip}{0.5em}

\setlist[itemize]{topsep=0.3em, itemsep=0.2em}
\setlist[enumerate]{topsep=0.3em, itemsep=0.2em}

\clubpenalty=10000
\widowpenalty=10000
\emergencystretch=1.5em

% =====================================================
% COMANDOS MATEMÁTICOS
% =====================================================
\newcommand{\abs}[1]{\left|#1\right|}
\newcommand{\norm}[1]{\left\lVert#1\right\rVert}
\newcommand{\vect}[1]{\mathbf{#1}}
\newcommand{\deriv}[2]{\frac{d#1}{d#2}}
\newcommand{\pderiv}[2]{\frac{\partial#1}{\partial#2}}

% =====================================================
% ENCABEZADOS Y PIES DE PÁGINA
% =====================================================
\pagestyle{fancy}
\fancyhf{}
\fancyhead[L]{\small\nouppercase{\leftmark}}
\fancyhead[R]{\small Lectocomprensión Jurídica}
\fancyfoot[C]{\thepage}
\renewcommand{\headrulewidth}{0.4pt}
\renewcommand{\footrulewidth}{0pt}

\fancypagestyle{plain}{
    \fancyhf{}
    \fancyfoot[C]{\thepage}
    \renewcommand{\headrulewidth}{0pt}
}

% =====================================================
% CAJAS ACADÉMICAS
% =====================================================
\newtcolorbox{definition}[1]{
    enhanced, breakable,
    title={Definición: #1},
    fonttitle=\bfseries,
    colback=blue!3,
    colframe=blue!50!black,
    boxrule=0.8pt, arc=2mm
}

\newtcolorbox{important}{
    enhanced, breakable,
    title={Importante},
    fonttitle=\bfseries,
    colback=yellow!8,
    colframe=orange!70!black,
    boxrule=0.8pt, arc=2mm
}

\newtcolorbox{example}[1]{
    enhanced, breakable,
    title={Ejemplo: #1},
    fonttitle=\bfseries,
    colback=green!4,
    colframe=green!40!black,
    boxrule=0.8pt, arc=2mm
}

\newtcolorbox{formula}[1]{
    enhanced, breakable,
    title={Regla / Fórmula: #1},
    fonttitle=\bfseries,
    colback=purple!4,
    colframe=purple!50!black,
    boxrule=0.8pt, arc=2mm
}

\newtcolorbox{note}{
    enhanced, breakable,
    title={Nota},
    fonttitle=\bfseries,
    colback=white,
    colframe=gray!50,
    boxrule=0.6pt, arc=2mm
}

% =====================================================
% DOCUMENTO
% =====================================================
\begin{document}

% ─────────────────────────────────────────────────────────
% PORTADA
% ─────────────────────────────────────────────────────────
\begin{titlepage}
\thispagestyle{empty}
\begin{center}

\vspace*{1cm}
{\large \textsc{\universidad}}

\vspace{0.2cm}
{\large \textsc{\facultad}}

\vspace{0.2cm}
{\large \carrera}

\vspace{3cm}
{\Huge \textbf{\materia}}

\vspace{1cm}
{\LARGE \titulo}

\vspace{0.8cm}
\textit{Comprender antes de traducir es el primer paso para transformar un texto ajeno en conocimiento propio.}

\vspace{1.2cm}
\rule{0.70\textwidth}{0.5pt}

\vspace{0.5cm}
{\large Apuntes de estudio integrados}

\vspace{0.5cm}
\rule{0.70\textwidth}{0.5pt}

\vfill
{\large \autor}

\vspace{0.3cm}
{\large \anio}

\end{center}

\vfill
\begin{flushleft}
\textit{Este material es un aporte personal al estudio de la materia. No pretende sustituir el material oficial de la cátedra ni los textos trabajados en clase.}
\end{flushleft}
\end{titlepage}

% ─────────────────────────────────────────────────────────
% ÍNDICE
% ─────────────────────────────────────────────────────────
\tableofcontents

% ─────────────────────────────────────────────────────────
% INTRODUCCIÓN GENERAL
% ─────────────────────────────────────────────────────────
\chapter*{Introducción general}
\addcontentsline{toc}{chapter}{Introducción general}
\markboth{Introducción general}{Introducción general}

Este manual reúne el contenido académico desarrollado a lo largo de la cursada de \textit{Lectocomprensión de Textos Jurídicos en Inglés}, organizado como un único cuerpo de conocimiento progresivo en lugar de una colección de clases independientes.

El recorrido se estructura en tres partes. La primera presenta los fundamentos metodológicos de la materia: qué se estudia, cómo se evalúa y, sobre todo, qué estrategias de lectura permiten abordar textos jurídicos en inglés sin depender de la traducción literal ni de la traducción automática. La segunda parte desarrolla la Unidad 1 del programa, dedicada al derecho constitucional comparado: primero el sistema de los Estados Unidos y su organización judicial, luego el sistema constitucional del Reino Unido con sus particularidades. La tercera parte desarrolla la Unidad 2, centrada en el derecho contractual del \textit{common law}: la formación del contrato, las figuras precontractuales, el análisis de un fallo seminal y las cláusulas estándar que aparecen en los contratos internacionales.

A lo largo de todo el manual, los contenidos lingüísticos y los contenidos jurídicos se presentan integrados, tal como ocurre en la práctica profesional: no se puede comprender un texto jurídico en inglés sin manejar las herramientas del idioma, ni tiene sentido aprender el idioma sin contenido jurídico que lo ancle.

% =====================================================
% PARTE I
% =====================================================
\part{Fundamentos metodológicos de la lectocomprensión jurídica}

% ─────────────────────────────────────────────────────────
\chapter{La materia: encuadre, objetivos y evaluación}
% ─────────────────────────────────────────────────────────

\section{¿Qué es Lectocomprensión de Textos Jurídicos en Inglés?}

La materia no tiene por objetivo enseñar inglés. Resulta imposible aprender a hablar un idioma en el término de un cuatrimestre, y ese no es su propósito. Lo que sí resulta posible, y constituye el eje central del trabajo, es que cada estudiante ---con el nivel de inglés que ya posea, mucho o poco--- desarrolle estrategias de lectura para abordar un texto jurídico en inglés, comprenderlo con mayor o menor profundidad, y explicar en español de qué se trata.

\begin{important}
El objetivo de la materia es que, con el conocimiento del idioma inglés que cada estudiante tenga y las estrategias de lectura trabajadas en la cursada, sea posible leer un texto jurídico en inglés y explicar en español su contenido general. No se trata de traducir palabra por palabra: se trata de comprender.
\end{important}

Esta distinción importa porque el ejercicio de cualquier profesión jurídica exige, cada vez más, manejarse con fuentes, doctrina, jurisprudencia y contratos redactados originalmente en inglés. Quien no desarrolla una estrategia de lectura eficiente queda dependiendo por completo de traducciones de terceros o de traductores automáticos que pueden inventar fallos, citas o normas inexistentes ---fenómeno conocido como ``alucinación'' de la inteligencia artificial, desarrollado más adelante.

La materia también enseña a tolerar la incertidumbre frente a un texto que no se comprende del todo, y a construir sentido igualmente, con lo que se sabe y con lo que se puede inferir: exactamente lo que se hace de manera naturalizada cuando se lee en la lengua materna.

\begin{example}{la lectura en lengua materna}
Al leer en español un autor de estilo complejo ---Borges, por ejemplo---, es habitual encontrar palabras cuyo significado exacto se desconoce. Sin embargo, la lectura continúa y, por el contexto, se comprende la idea general del texto, aun sin conocer el significado preciso de cada término. Esta misma estrategia es la que la materia busca trasladar a la lectura de textos jurídicos en inglés.
\end{example}

\begin{example}{el texto en rumano}
En una capacitación docente, se entregó a un grupo de profesores un texto escrito en rumano ---idioma que ninguno conocía--- con la consigna de identificar la idea principal en cinco minutos. El primer impulso fue recurrir a la traducción automática; sin embargo, apelando a estrategias propias de quien enfrenta un idioma extranjero, el grupo logró aproximarse al sentido general del texto. Esta experiencia ilustra la situación del estudiantado frente a un texto jurídico en inglés: con estrategia, la comprensión es posible aun sin dominar el idioma.
\end{example}

\section{Ubicación en el plan de estudios y relación con el derecho comparado}

La materia se ubica hacia el final de la carrera porque requiere que el estudiantado cuente ya con una formación jurídica sólida en el derecho argentino. Esa formación previa es una condición necesaria para poder comparar el sistema argentino con el sistema jurídico extranjero del \textit{common law}.

Sin un conocimiento previo de nociones como constitución rígida o flexible, o del funcionamiento del proceso civil argentino, resulta imposible explicar las diferencias con el proceso anglosajón. El vocabulario jurídico constituye un lenguaje de especialidad tanto en español como en inglés, y esta especificidad terminológica es independiente del dominio general del idioma. Términos como \textit{usucapión}, \textit{usufructo} o \textit{servidumbre} resultan ajenos a quien no tiene formación jurídica, aun siendo hablante nativo de español; el mismo fenómeno se reproduce en inglés respecto del vocabulario técnico del derecho.

A medida que avanza el programa, las semejanzas entre el sistema argentino (de tradición continental) y el sistema anglosajón (\textit{common law}) tienden a disminuir: son mayores en la primera unidad, referida al derecho constitucional, y se reducen progresivamente al abordar contratos y derecho procesal, en la medida en que se trata de dos tradiciones jurídicas de base distinta.

\begin{table}[H]
\centering
\caption{Contraste entre el sistema continental y el sistema anglosajón}
\begin{tabularx}{\textwidth}{>{\raggedright\arraybackslash}p{0.22\textwidth} >{\raggedright\arraybackslash}X >{\raggedright\arraybackslash}X}
\toprule
\textbf{Aspecto} & \textbf{Sistema continental (Argentina)} & \textbf{Sistema anglosajón (\textit{common law})} \\
\midrule
Origen y base & Derecho romano-germánico; normas codificadas & Derecho consuetudinario (\textit{customary law}); precedentes judiciales \\
Peso de la palabra & Predomina lo escrito (escritos, expedientes) & Predomina lo oral (audiencias, testimonios, alegatos) \\
Proceso judicial & Mayormente escrito, aunque con expansión del juicio oral & Mayormente oral; juicio por jurados \\
Fuente del derecho & Constitución, tratados, leyes y decretos & Precedentes judiciales (\textit{case law}) y legislación \\
\bottomrule
\end{tabularx}
\end{table}

\section{El uso de la inteligencia artificial y los traductores automáticos}
\label{sec:ia}

La cátedra reconoce abiertamente que utiliza la inteligencia artificial como herramienta de trabajo, en particular para la preparación de clases. No se sostiene una postura de rechazo tecnológico. Sin embargo, se advierte sobre dos riesgos concretos.

El primero es el de las \textbf{``alucinaciones''}: al solicitar a un sistema de inteligencia artificial que elabore un informe con jurisprudencia, o que sugiera fuentes sobre un tema jurídico, el sistema puede ofrecer resultados aparentemente correctos ---enlaces, fallos, sitios--- que en la práctica no existen. El sistema puede inventar fallos, citas o fuentes inexistentes con apariencia de verosimilitud. Esto exige una revisión crítica constante de toda información obtenida por este medio.

\begin{example}{las ``alucinaciones'' de la inteligencia artificial}
Al solicitar a un sistema de inteligencia artificial que sugiera jurisprudencia sobre determinado tema, el sistema puede ofrecer nombres de fallos, caratulas y números de expediente que tienen apariencia de verosimilitud pero que en la práctica no existen. Solo quien tiene formación suficiente en el área puede detectar el error.
\end{example}

El segundo riesgo es el de la \textbf{dependencia}: el uso indiscriminado de la inteligencia artificial para resolver un texto equivale a pedir comida a domicilio mientras se pretende aprender a cocinar. El resultado inmediato se obtiene, pero no se produce aprendizaje alguno. En la materia se solicita frenar el impulso de pasar todo el texto por el traductor automático, tanto porque no podrá utilizarse el día del examen como porque dicha práctica impide el desarrollo del criterio propio.

\begin{example}{el cálculo matemático sin herramientas}
Quien no aprendió a calcular un porcentaje de manera autónoma no cuenta con criterio propio para advertir si el resultado que arroja una calculadora es correcto. Solo quien desarrolló la habilidad de base puede detectar un error. Lo mismo ocurre con la inteligencia artificial: solo quien tiene formación suficiente puede advertir cuándo su resultado es erróneo.
\end{example}

\begin{important}
La inteligencia artificial y los traductores automáticos son herramientas que llegaron para quedarse y que no corresponde ignorar. Pero su utilización sin criterio propio conlleva el riesgo de una dependencia que anula la capacidad de resolver problemas de manera autónoma. La materia busca construir ese criterio propio.
\end{important}

\section{Organización de la cursada}

La modalidad habitual es presencial. La virtualidad se reserva para situaciones excepcionales ---paros docentes, movilizaciones, cierres imprevistos--- y no constituye la modalidad habitual de trabajo.

Se requiere un 75\% de asistencia para la regularidad. Las inasistencias justificadas ---por cuestiones médicas o coincidencia con audiencias de prácticas profesionales--- deben acreditarse en el momento en que se producen, no de forma extemporánea.

El material de cátedra se comparte a través de la plataforma Classroom, accesible obligatoriamente con el correo institucional. El programa de la materia puede entenderse como un \textbf{contrato pedagógico}: la cátedra asume una prestación (el dictado y la orientación del aprendizaje) y el estudiantado asume una contraprestación (el control de la propia asistencia, la concurrencia con el material, la resolución de actividades).

\subsection{Diccionarios y fuentes de referencia}

Se requieren dos tipos de diccionario: uno general (Larousse, Collins o Webster's) y uno específicamente jurídico bilingüe. El diccionario general no incluye terminología técnica del derecho; de ahí la necesidad del diccionario jurídico.

Se recomienda especialmente el diccionario bilingüe jurídico de Cabanellas (tomo inglés-español), el único tomo necesario dado que en la materia no se producirá ningún texto en inglés. Una característica de ese diccionario es que, en lugar de ofrecer un término equivalente preciso, suele dar la definición completa del concepto, lo cual resulta especialmente útil para comprender institutos del sistema anglosajón que no existen o difieren del sistema argentino.

\begin{important}
Para el primer parcial es obligatorio contar con diccionario en formato papel. No está permitido el uso de dispositivos electrónicos durante el examen. Quien no disponga de diccionario propio debe solicitarlo prestado o retirarlo de la biblioteca con antelación.
\end{important}

\section{Sistema de evaluación}
\label{sec:evaluacion}

\subsection{Primer parcial}

El primer parcial es individual, presencial y escrito. Se toma un día de la semana con tres horas disponibles ---aunque está pensado para completarse en aproximadamente hora y media o dos horas--- y requiere el uso de diccionario en papel. Consiste en la lectura de un texto en inglés sobre un tema trabajado en clase y en la resolución de consignas redactadas en español. No se exige ninguna producción escrita en inglés.

\begin{important}
La redacción en español constituye una condición central de aprobación del parcial, independientemente de la comprensión efectiva del texto en inglés. Una comprensión correcta que no pueda expresarse con claridad en español puede determinar la desaprobación del examen. El lenguaje escrito es especialmente relevante para el ejercicio de la disciplina jurídica, donde la forma en que se expresa la comprensión de un texto resulta determinante.
\end{important}

El parcial suele tomarse más tarde que en otras materias cuatrimestrales ---hacia fines de octubre--- porque el objetivo de la cátedra es la incorporación de estrategias de lectura, un proceso que requiere mayor tiempo de maduración que la acumulación de contenido teórico.

El recuperatorio, destinado a quienes no aprueben el primer parcial, sigue un formato similar aunque algo más breve, y se toma aproximadamente una semana después de conocidas las notas.

\subsection{Trabajo grupal de investigación y exposición oral}

Superado el primer parcial, se habilita una segunda instancia evaluativa: un trabajo grupal, en general de cuatro integrantes, con una etapa escrita y una etapa oral. En la etapa escrita, cada grupo elige un tema jurídico no trabajado previamente en clase y elabora un trabajo de investigación a partir de varias fuentes en inglés. En la etapa oral, cada grupo realiza una exposición de entre 15 y 20 minutos.

\begin{example}{la especialización notarial comparada}
En un cuatrimestre anterior, un grupo eligió como tema la especialización notarial, comparando la figura del escribano público argentino con la figura equivalente en Estados Unidos, que presenta numerosas diferencias respecto de la argentina. Este ejemplo se destaca como modelo de elección de un tema original.
\end{example}

La calificación de ambas instancias es binaria: aprobado o desaprobado, sin nota numérica. Las exposiciones orales grupales se concentran en las últimas clases del cuatrimestre.

% ─────────────────────────────────────────────────────────
\chapter{Estrategias de lectocomprensión}
\label{cap:estrategias}
% ─────────────────────────────────────────────────────────

Este capítulo reúne las herramientas metodológicas centrales de la materia. No se presentan como reglas a memorizar, sino como procedimientos que se van automatizando con la práctica. La pregunta de fondo que orienta todas estas estrategias es siempre la misma: ¿cómo puedo extraer el sentido general de este texto sin entender cada una de sus palabras?

\section{La macroestructura del texto}
\label{sec:macroestructura}

La primera estrategia de lectura consiste en observar el texto antes de leerlo en profundidad: mirar su estructura visual antes de descender al detalle léxico.

\begin{definition}{macroestructura del texto}
La macroestructura de un texto es el conjunto de elementos paratextuales ---títulos, subtítulos, viñetas, colores, subrayados, tipografía y fuente de publicación--- que pueden observarse antes de una lectura en profundidad. Permiten anticipar el contenido general y organizar la información previamente a descender al detalle léxico.
\end{definition}

Esta estrategia se aplica de manera automática en español ---por ejemplo, al identificar en un manual de derecho la estructura típica de naturaleza jurídica, elementos, definición y clasificación--- y simplemente debe hacerse consciente al leer en inglés.

Un elemento importante de la macroestructura es la \textbf{fuente del texto}: no es equivalente obtener información sobre la estructura del Estado de los Estados Unidos de una fuente como Wikipedia que de un sitio oficial como \textit{usa.gov}. La fuente anticipan el registro, la confiabilidad y, a veces, el contenido del texto.

\begin{note}
La presencia de subtítulos en un texto permite anticipar que cada apartado ofrecerá una definición o conceptualización del término correspondiente, antes incluso de comenzar la lectura detallada. Cuando un texto carece de subtítulos ---como el texto ``What is the UK Constitution?''---, la primera lectura debe apoyarse en otros indicios: el título, la fuente y los términos reconocibles.
\end{note}

\section{Términos transparentes y cognados}

\begin{definition}{cognado o término transparente}
Un cognado o término transparente es una palabra que, al leerse o escucharse en otra lengua, permite reconocer de forma casi inmediata la idea que representa en la lengua propia, generalmente por compartir un origen etimológico común. El español, en tanto lengua romance derivada del latín, y el inglés, lengua de origen sajón, tienen orígenes distintos; sin embargo, el inglés también incorporó términos provenientes del latín a lo largo de su historia, lo que genera estos puntos de contacto.
\end{definition}

A partir de un texto sobre la organización del Estado de los Estados Unidos, pueden identificarse numerosos cognados: \textit{Constitution/Constitución}, \textit{government/gobierno}, \textit{legislative/legislativo}, \textit{executive/ejecutivo}, \textit{judicial/judicial}, \textit{president/presidente}, \textit{cabinet/gabinete}, \textit{Congress/Congreso}, \textit{Supreme Court/Corte Suprema}, \textit{federal/federal}, entre otros.

El valor pedagógico de esta estrategia es especialmente relevante para quienes tienen escaso dominio del inglés: permite reducir significativamente la sensación de desconocimiento total frente a un texto. Aun sin poder describir la totalidad de su contenido, el reconocimiento de los cognados permite ubicarse temáticamente.

\begin{note}
No toda palabra reconocible por contexto jurídico es necesariamente un cognado en sentido estricto. Los términos \textit{branch} (poder del Estado) y \textit{law} (ley) pueden reconocerse por asociación conceptual previa, no por semejanza fonética o gráfica directa con su equivalente en español. Esto no los hace menos útiles para la comprensión, pero sí impide clasificarlos como cognados.
\end{note}

\subsection{Palabras largas en inglés}

El inglés tiende a tener palabras cortas, de dos o tres sílabas. Cuando aparecen palabras muy largas, generalmente contienen prefijos, sufijos, o son de origen latino. El origen común de estos términos con el español los hace, en muchos casos, reconocibles. Esto explica por qué las palabras técnicas del derecho ---\textit{unconstitutional}, \textit{jurisdiction}, \textit{legislative}--- resultan más accesibles que el vocabulario cotidiano.

\section{Falsos cognados (\textit{false friends})}
\label{sec:falsos-cognados}

\begin{definition}{falso cognado (\textit{false friend})}
Un falso cognado o falso amigo (\textit{false friend}) es una palabra que, por su semejanza formal con un término de otra lengua, parece significar una determinada cosa cuando en realidad tiene un significado distinto en el contexto dado. Se detecta cuando, durante la lectura, algo ``no cierra'' o genera una sensación de extrañeza respecto del sentido de la oración.
\end{definition}

No existe una lista de falsos cognados que pueda aprenderse de memoria como los verbos irregulares. La forma de identificarlos es la siguiente: cuando se está leyendo un texto e interpretando su contenido, surge la pregunta interna ``¿qué tiene que ver esto aquí?''. Cuando una palabra no cierra y la lectura pierde coherencia, es una señal de que puede tratarse de un falso cognado. La solución es buscar esa palabra en el diccionario.

El conocimiento jurídico previo es fundamental para detectar los falsos cognados: es quien conoce el derecho quien puede advertir que una traducción literal produce un resultado jurídicamente absurdo.

Los principales falsos cognados jurídicos aparecen desarrollados en su contexto a lo largo de este manual. El siguiente cuadro ofrece una síntesis de referencia:

\begin{longtable}{
    >{\raggedright\arraybackslash}p{0.24\textwidth}
    >{\raggedright\arraybackslash}p{0.28\textwidth}
    >{\raggedright\arraybackslash}p{0.37\textwidth}
}
\caption{Falsos cognados jurídicos identificados en la cursada} \\
\toprule
\textbf{Término en inglés} & \textbf{Falso sentido intuitivo} & \textbf{Significado correcto en contexto} \\
\midrule
\endfirsthead
\toprule
\textbf{Término en inglés} & \textbf{Falso sentido intuitivo} & \textbf{Significado correcto en contexto} \\
\midrule
\endhead
\textit{government} & Gobierno (poder ejecutivo) & Estado / organización estatal en su conjunto \\
\textit{justice(s)} & Justicia (abstracto) & Ministro/s de la Corte Suprema \\
\textit{jurisprudence} & Jurisprudencia (fallos judiciales) & Filosofía del derecho; doctrina \\
\textit{statute} & Estatuto (norma interna de organización) & Ley formal sancionada por el Parlamento \\
\textit{convention} & Convención internacional; tratado & Costumbre o práctica institucional no escrita \\
\textit{court} & Corte (de cualquier nivel) & Tribunal (genérico); ``Corte'' solo para la Suprema \\
\textit{sentence} & Sentencia (judicial en general) & Pena impuesta a quien fue declarado culpable \\
\textit{jurisdiction} & Jurisdicción (en sentido abstracto) & Competencia (qué juez entiende en cada causa) \\
\textit{consideration} & Consideración & Contraprestación (en derecho contractual) \\
\textit{evidence} & Evidencia & Prueba; medios de prueba \\
\textit{checks and balances} & Cheques y balances (literal) & Sistema de frenos y contrapesos \\
\textit{support} (en contexto) & Soportar (tolerar) & Apoyar; asistir \\
\textit{term} (en contexto político) & Término (palabra) & Mandato; período de ejercicio de un cargo \\
\textit{bargain} (en contexto jurídico) & Regatear & Negociación precontractual; tratativas \\
\textit{will} (en contexto sucesorio) & Voluntad (abstracta) & Testamento \\
\bottomrule
\end{longtable}

\section{Análisis morfológico: prefijos, sufijos y palabras compuestas}

Muchas palabras en inglés pueden descomponerse en partes que permiten inferir su significado sin necesidad del diccionario.

\begin{formula}{Prefijos frecuentes en inglés jurídico}
\begin{itemize}
    \item \textbf{un-} (no, in-): \textit{unconstitutional} (inconstitucional), \textit{unwritten} (no escrito), \textit{unchallengeable} (incuestionable).
    \item \textbf{en-} (dar fuerza a): \textit{enforce} (hacer cumplir; literalmente ``darle fuerza'').
    \item \textbf{over-} (sobre, exceder): \textit{overturn} (anular; literalmente ``dar vuelta sobre'').
    \item \textbf{dis-} (negación): \textit{discharge} (extinción del contrato; en contexto general, liberar de una obligación).
\end{itemize}
\end{formula}

\begin{formula}{Sufijos frecuentes en inglés jurídico}
\begin{itemize}
    \item \textbf{-able} (que puede ser): \textit{enforceable} (exigible), \textit{unchallengeable} (incuestionable).
    \item \textbf{-ness} (sustantivador abstracto): \textit{willingness} (voluntad), a partir de \textit{willing}.
    \item \textbf{-tion} (nominalizador de verbos): \textit{termination}, \textit{jurisdiction}, \textit{novation}.
    \item \textbf{-or / -ee} (sujeto activo / sujeto receptor en negocios jurídicos): \textit{offeror/offeree}, \textit{employer/employee}, \textit{franchisor/franchisee}, \textit{promisor/promisee}, \textit{grantor/grantee}. La terminación \textbf{-or} indica quien origina el acto; la terminación \textbf{-ee} indica quien lo recibe o resulta beneficiario. Esta distinción, que en inglés se marca con un sufijo, no siempre tiene traducción natural en español: \textit{offeree}, por ejemplo, no tiene un equivalente único en nuestra lengua.
\end{itemize}
\end{formula}

Las \textbf{palabras compuestas} siguen una lógica similar: \textit{landmark} (hito; \textit{land} = tierra + \textit{mark} = marca) designa en derecho las leyes o fallos que marcaron un punto de inflexión histórico (\textit{landmark laws}, \textit{leading cases}).

\section{Colocaciones léxicas (\textit{collocations})}

\begin{definition}{colocación (\textit{collocation})}
Una colocación es una combinación tradicional de palabras que aparecen frecuentemente juntas en una lengua. Las colocaciones no se deducen del diccionario: requieren exposición repetida a textos auténticos y la construcción de un repertorio de combinaciones típicas del campo de especialidad.
\end{definition}

Las colocaciones son especialmente importantes en el lenguaje jurídico, donde la combinación incorrecta puede cambiar el sentido de una cláusula.

\begin{table}[H]
\centering
\caption{Colocaciones jurídicas frecuentes en inglés}
\begin{tabularx}{\textwidth}{>{\raggedright\arraybackslash}p{0.33\textwidth} >{\raggedright\arraybackslash}X}
\toprule
\textbf{Expresión en inglés} & \textbf{Equivalente en español} \\
\midrule
To pass a law & Sancionar / promulgar una ley \\
To celebrate / enter into a contract & Celebrar / formalizar un contrato \\
To sign / execute an agreement & Firmar / celebrar un acuerdo \\
To enforce a law / a contract & Hacer cumplir una ley / un contrato \\
To file / lodge a claim & Interponer una demanda \\
To bring an action & Iniciar una acción judicial \\
To hand down a decision & Dictar sentencia \\
To draft a contract & Redactar un contrato \\
To comply with & Cumplir con (una obligación) \\
\bottomrule
\end{tabularx}
\end{table}

\begin{note}
La expresión ``to celebrate a contract'' en inglés jurídico no significa festejar: significa formalizar o concluir un contrato. ``To enter into a contract'' no alude a un ingreso físico sino al momento en que las partes formalizan el acuerdo. Estas colocaciones son propias del campo especializado y no pueden deducirse de reglas generales.
\end{note}

\section{Verbos modales en el lenguaje jurídico}
\label{sec:modales}

Los verbos modales son auxiliares que, por sí solos, no tienen significado léxico; necesitan ir acompañados de un verbo principal. En el lenguaje jurídico expresan, de manera precisa, el \textbf{deber ser}: obligación, facultad, permiso, recomendación o posibilidad. Esta distinción resulta central en contratos y normas, donde la diferencia entre una obligación y una facultad puede cambiar radicalmente el alcance de una cláusula.

\begin{table}[H]
\centering
\caption{Verbos modales en el lenguaje jurídico en inglés}
\begin{tabularx}{\textwidth}{>{\raggedright\arraybackslash}p{0.13\textwidth} >{\raggedright\arraybackslash}p{0.32\textwidth} >{\raggedright\arraybackslash}X}
\toprule
\textbf{Modal} & \textbf{Idea que expresa} & \textbf{Ejemplo} \\
\midrule
\textbf{must} & Obligación ineludible. La más fuerte; emana de una ley o contrato. & \textit{There must be a quorum of six} = se necesitan seis jueces; requisito ineludible. \\
\textbf{shall} & La obligación por excelencia del lenguaje jurídico. Propia de contratos, legislación y documentos religiosos. Poco frecuente en lenguaje cotidiano. & \textit{The seller shall deliver the goods} = el vendedor deberá entregar los bienes. \\
\textbf{may} & Potestad o facultad (no obligación); también, permiso. & \textit{The court may decide} = el tribunal tiene la facultad de decidir. \\
\textbf{might} & Posibilidad más remota o incierta. & \textit{The claim might succeed} = el reclamo podría prosperar (con menor probabilidad). \\
\textbf{should} & Recomendación o sugerencia; no obligación. & \textit{You should see a professional} = debería consultar a un profesional. \\
\textbf{would / could} & Condicional; pedidos educados. & \textit{Could you please confirm?} \\
\bottomrule
\end{tabularx}
\end{table}

La distinción entre \textit{must} y \textit{shall} es particularmente relevante. Ambos expresan obligación, pero \textit{shall} connota una autoridad superior ---la ley, el contrato, la norma--- que impone ese deber. En el lenguaje jurídico español, el uso del futuro con sentido de obligación es el equivalente más frecuente: ``las partes deberán acordar'', ``el comprador deberá abonar el precio''.

\begin{example}{Shall en contratos}
``The purchaser shall pay the price'' (el comprador deberá abonar el precio). ``The seller shall deliver the goods'' (el vendedor deberá entregar los bienes). Estas obligaciones emanan de lo pactado entre las partes, que a su vez encuentra fundamento en la norma que reconoce la fuerza vinculante de los contratos con objeto lícito.
\end{example}

\section{La estrategia de fraccionamiento de oraciones complejas}

Cuando una oración es compleja o genera dificultad, una estrategia efectiva consiste en extraer las partes secundarias ---habitualmente aquellas que van entre comas, como las proposiciones relativas--- y leer primero la oración simplificada. Luego se reintegran las partes extraídas.

\begin{example}{fraccionamiento de una oración compleja}
Oración original: \textit{``The Justices of the Supreme Court, who are nominated by the president and confirmed by the Senate, can overturn unconstitutional laws.''}

Versión simplificada: \textit{``The Justices of the Supreme Court can overturn unconstitutional laws.''}

Luego se agrega la información entre comas: los jueces son nominados por el presidente y confirmados por el Senado. El sentido global resulta más accesible trabajando por partes.
\end{example}

Esta técnica es especialmente útil con las \textbf{estructuras nominales encadenadas}, características del inglés: ``\textit{a two-year term}'' (un mandato de dos años) o ``\textit{landmark constitutional laws}'' (leyes constitucionales fundamentales). El inglés permite encadenar sustantivos y adjetivos sin preposiciones; al traducir, es necesario agregar esas preposiciones y reestructurar la frase.

\section{La importancia del contexto y cuándo usar el diccionario}

Un mismo término puede tener significados completamente distintos según el contexto en que aparezca. Esta es quizás la lección metodológica más importante de la materia: antes de buscar en el diccionario, leer el contexto.

\begin{itemize}
    \item \textit{Government} puede significar ``gobierno'' (poder ejecutivo) en algunos contextos, y ``Estado'' en otros.
    \item \textit{State} puede referirse al Estado federal (\textit{the State}) o a un estado particular (\textit{the states}); la diferencia entre singular y plural, o entre mayúscula y minúscula, es relevante.
    \item \textit{Law} puede significar una norma en particular, el sistema jurídico en su conjunto, un área del derecho, o coloquialmente la policía o la profesión legal (véase el Capítulo~\ref{cap:poder-judicial}).
    \item \textit{Court} puede referirse a cualquier nivel de tribunal, no necesariamente a la Corte Suprema (véase Sección~\ref{sec:court}).
    \item \textit{Term} puede significar ``término'' (palabra) o ``mandato'' (período de ejercicio de un cargo).
    \item \textit{Will} puede significar ``voluntad'' o ``testamento'' (en contexto sucesorio).
\end{itemize}

La regla general para usar el diccionario es: consultarlo cuando la incomprensión de un término impide entender la idea central de una oración o párrafo. No todo desconocimiento requiere consultar el diccionario; las palabras del lenguaje cotidiano generalmente pueden deducirse del contexto. Los términos técnicos del derecho, en cambio, suelen requerir verificación, dado que su significado específico puede diferir radicalmente del sentido coloquial.

% ─────────────────────────────────────────────────────────
% CUADRO CORNELL — PARTE I
% ─────────────────────────────────────────────────────────
\chapter*{Cuadro Cornell: Fundamentos metodológicos}
\addcontentsline{toc}{chapter}{Cuadro Cornell: Fundamentos metodológicos}
\markboth{Cuadro Cornell: Fundamentos metodológicos}{Cuadro Cornell: Fundamentos metodológicos}

\begin{longtable}{
    >{\raggedright\arraybackslash}p{0.30\textwidth}
    >{\raggedright\arraybackslash}p{0.60\textwidth}
}
\toprule
\textbf{Preguntas / palabras clave} & \textbf{Notas / respuestas} \\
\midrule
\endfirsthead
\toprule
\textbf{Preguntas / palabras clave} & \textbf{Notas / respuestas} \\
\midrule
\endhead
\textbf{¿Cuál es el objetivo de la materia?} &
Desarrollar estrategias de lectura para abordar textos jurídicos en inglés y poder explicar en español de qué tratan, sin necesidad de traducciones literal ni automáticas. No es un curso de idioma inglés. \\

\textbf{¿Qué es la macroestructura?} &
El conjunto de elementos paratextuales (títulos, subtítulos, viñetas, fuente del texto) que permiten anticipar el contenido antes de una lectura en profundidad. Es la primera estrategia de lectura. \\

\textbf{¿Qué es un cognado y cómo se reconoce?} &
Una palabra que, por compartir origen etimológico con su equivalente en español (generalmente latín), se reconoce con facilidad. Ejemplos: \textit{Constitution}, \textit{federal}, \textit{executive}. \\

\textbf{¿Cómo se detecta un falso cognado?} &
Cuando durante la lectura algo ``no cierra'': la oración pierde coherencia o el resultado jurídico es absurdo. Se verifica con el diccionario. No existe una lista para memorizar. \\

\textbf{¿Qué expresan los sufijos -or y -ee?} &
\textit{-or} identifica al sujeto activo que origina el acto (offeror, employer); \textit{-ee} identifica al receptor o beneficiario (offeree, employee). Provienen del francés y se reproducen en la mayoría de los negocios jurídicos en inglés. \\

\textbf{¿Qué diferencia hay entre \textit{must} y \textit{shall}?} &
Ambos expresan obligación. \textit{Must}: obligación emanada de la ley o el contrato. \textit{Shall}: la obligación por excelencia del lenguaje jurídico, con connotación de autoridad superior; poco frecuente en el lenguaje cotidiano. \\

\textbf{¿Qué diferencia hay entre \textit{may} y \textit{might}?} &
\textit{May}: potestad o facultad (probabilidad alta o capacidad legal). \textit{Might}: posibilidad más remota (probabilidad menor). \\

\textbf{¿Qué es una colocación?} &
Una combinación tradicional de palabras que aparecen juntas en un campo especializado. No se deduce del diccionario. Ejemplos: \textit{to pass a law} (sancionar una ley), \textit{to enter into a contract} (celebrar un contrato). \\

\textbf{¿Cuándo conviene usar el diccionario?} &
Cuando la incomprensión de un término impide entender la idea central de una oración o párrafo, especialmente con términos técnicos cuyo significado específico difiere del sentido coloquial. \\

\textbf{¿Qué riesgos presenta la inteligencia artificial?} &
Las ``alucinaciones'' (fallos, citas y fuentes inexistentes con apariencia verosímil) y la dependencia (quien no desarrolló la habilidad de base no puede detectar el error). Se admite como herramienta de apoyo; se rechaza como sustituto del criterio propio. \\

\midrule
\multicolumn{2}{p{0.90\textwidth}}{
\textbf{Resumen:} La materia no enseña inglés: enseña a leer estratégicamente un texto jurídico en un idioma que no se domina del todo. Las herramientas centrales son: la lectura de la macroestructura como primer acercamiento; el reconocimiento de cognados para reducir la sensación de desconocimiento; la detección de falsos cognados cuando el sentido de una oración genera extrañeza; el análisis morfológico de prefijos y sufijos para inferir significados; el conocimiento de colocaciones propias del campo jurídico; y la interpretación precisa de los verbos modales, que en el lenguaje jurídico expresan con exactitud la naturaleza de cada obligación o facultad. A estas herramientas se suma una postura crítica frente a la inteligencia artificial: útil como apoyo, peligrosa como sustituto.
} \\
\bottomrule
\end{longtable}

% =====================================================
% PARTE II
% =====================================================
\part{Unidad 1: Derecho constitucional comparado}

% ─────────────────────────────────────────────────────────
\chapter{La organización del Estado de los Estados Unidos}
\label{cap:eeuu}
% ─────────────────────────────────────────────────────────

Los textos trabajados en la primera unidad del programa provienen de fuentes oficiales del gobierno de los Estados Unidos. Su lenguaje es intencionalmente accesible: están pensados para que cualquier ciudadano, sin formación jurídica, pueda comprender la organización del Estado. Esto explica el uso de verbos simples (\textit{makes}, \textit{carries out}, \textit{evaluates}) en lugar de la terminología técnica más precisa que utilizaría un abogado.

\section{El término \textit{government}: un falso cognado fundamental}
\label{sec:government}

El primer elemento a resolver antes de avanzar en la lectura de cualquier texto sobre la organización estatal estadounidense es la palabra \textit{government}.

La tendencia natural es traducirla como ``gobierno'' en el sentido del poder ejecutivo: el presidente y los ministros. Sin embargo, en estos textos \textit{government} refiere a toda la estructura del Estado ---ejecutivo, legislativo y judicial--- y resulta equivalente, en este contexto, al concepto argentino de ``Estado''. La prueba es sencilla: el texto afirma que el \textit{government} se divide en tres poderes, lo cual resultaría absurdo si ``gobierno'' significara solo el ejecutivo.

\begin{important}
En inglés, \textit{government} se utiliza habitualmente para referirse a la estructura estatal completa. En español, ``gobierno'' se reserva generalmente al poder ejecutivo. Esta diferencia de uso hace que \textit{government} funcione como un falso cognado en los contextos constitucionales, donde debe traducirse como ``Estado'' u ``organización estatal''.
\end{important}

\section{La división de poderes y el sistema de frenos y contrapesos}

La Constitución de los Estados Unidos organiza el gobierno federal en tres poderes o ramas (\textit{branches}). El objetivo de esta división es que ningún individuo o grupo concentre demasiado poder.

\begin{note}
La traducción literal de \textit{branch} es ``rama'' (de árbol). En el contexto constitucional, sin embargo, resulta más preciso hablar de ``poderes'' que de ``ramas''. Aunque ``rama'' no es técnicamente incorrecto, ``poder'' es el término que refleja mejor la realidad jurídica de estas divisiones.
\end{note}

\begin{table}[H]
\centering
\caption{Los tres poderes del gobierno federal y sus funciones}
\begin{tabularx}{\textwidth}{>{\raggedright\arraybackslash}p{0.26\textwidth} >{\raggedright\arraybackslash}X}
\toprule
\textbf{Poder} & \textbf{Función según el texto} \\
\midrule
Legislativo (\textit{legislative branch}) & Elabora, sanciona y redacta las leyes (\textit{makes laws}). \\
Ejecutivo (\textit{executive branch}) & Lleva a cabo y hace cumplir las leyes (\textit{carries out and enforces laws}). \\
Judicial (\textit{judicial branch}) & Evalúa e interpreta las leyes; controla su constitucionalidad (\textit{evaluates laws}). \\
\bottomrule
\end{tabularx}
\end{table}

El sistema completo se denomina \textit{system of checks and balances} ---sistema de frenos y contrapesos--- y establece que cada poder puede controlar, limitar y equilibrar el poder de los otros. La traducción literal de ``checks and balances'' como ``cheques y balances'' no tiene sentido: \textit{check} en este contexto significa ``control'' o ``revisión'', y \textit{balance} significa ``equilibrio'', no ``balance contable''.

\section{El verbo \textit{enforce}: hacer cumplir la ley}

El verbo \textit{enforce} es sumamente jurídico. Compuesto por el prefijo \textit{en-} y la palabra \textit{force} (fuerza), literalmente significa ``darle fuerza''. En contexto legal, su significado es hacer cumplir o poner en práctica una norma. Esta es la función específica del poder ejecutivo en el texto trabajado.

\section{El poder legislativo: el Congreso estadounidense}
\label{sec:congreso}

El Congreso está compuesto por dos cámaras, lo cual configura un sistema bicameral similar al argentino.

\begin{note}
La traducción literal de \textit{House} es ``casa''. En el contexto del Congreso estadounidense, se traduce como ``Cámara'': \textit{House of Representatives} es la ``Cámara de Representantes''. En español también se la denomina ``Cámara de Diputados''.
\end{note}

\begin{table}[H]
\centering
\caption{Composición del Congreso de los Estados Unidos}
\begin{tabularx}{\textwidth}{>{\raggedright\arraybackslash}p{0.20\textwidth} >{\raggedright\arraybackslash}X >{\raggedright\arraybackslash}X >{\raggedright\arraybackslash}p{0.20\textwidth}}
\toprule
\textbf{Cámara} & \textbf{Composición} & \textbf{Mandato} & \textbf{Reelección} \\
\midrule
Senado & 100 senadores: 2 por cada estado, independientemente de su población. & 6 años & Ilimitada \\
Cámara de Representantes & 435 representantes, distribuidos entre los 50 estados en proporción a su población. & 2 años & Ilimitada \\
\bottomrule
\end{tabularx}
\end{table}

La composición del Senado refleja una decisión política deliberada: al asignar dos senadores a cada estado con independencia de su población, Wyoming (escasamente poblado) tiene la misma representación que California (el estado más poblado). La Cámara de Representantes, en cambio, introduce la proporcionalidad poblacional.

El texto menciona también la existencia de delegados no votantes (\textit{nonvoting delegates}) que representan al Distrito de Columbia y a los territorios estadounidenses. El Distrito de Columbia es similar a la Ciudad Autónoma de Buenos Aires: una ciudad capital con estatus especial que no forma parte de ningún estado.

El término \textit{term} en este contexto político no refiere a una palabra, sino al mandato o período de ejercicio de un cargo. ``A two-year term'' significa ``un mandato de dos años''. Esta es una de las formas en que el inglés encadena sustantivos y adjetivos sin necesidad de preposiciones: ``a two-year term'' es, literalmente, un término-de-dos-años.

\section{El poder ejecutivo}

El poder ejecutivo se encarga de ejecutar y hacer cumplir las leyes que sanciona el legislativo.

\subsection{Presidente, vicepresidente y gabinete}

El Presidente es el jefe del Estado, líder del gobierno federal y comandante en jefe de las Fuerzas Armadas. Cumple un mandato de \textbf{4 años} y puede ser elegido no más de dos veces (máximo 8 años de gobierno). Este límite fue establecido constitucionalmente después de que Franklin D. Roosevelt se reeligiera cuatro veces.

El Vicepresidente asiste al Presidente y lo reemplaza automáticamente si este no puede ejercer sus funciones por cualquier motivo. A diferencia del Presidente, el Vicepresidente puede ser reelecto un número ilimitado de veces.

El Gabinete (\textit{Cabinet}) está integrado por funcionarios de alto rango (\textit{high ranking government officials}) designados por el Presidente. Sus miembros deben ser confirmados por el Senado mediante mayoría simple (51 votos de 100). Este es uno de los mecanismos de freno del legislativo sobre el ejecutivo.

\subsection{Facultades del presidente}

El Presidente tiene, entre otras, las siguientes facultades: vetar leyes del Congreso; nominar a los jefes de agencias federales; y nominar a los jueces federales, incluidos los de la Corte Suprema.

\section{El poder judicial}
\label{sec:judicial-eeuu}

El poder judicial es responsable de interpretar las leyes y controlar su constitucionalidad. La Corte Suprema puede declarar una ley inconstitucional, anulando (\textit{overturning}) su vigencia.

El término \textit{overturn} es una palabra compuesta: \textit{turn} (girar, dar vuelta) más \textit{over} (sobre). En el contexto constitucional, su sentido es ``declarar inconstitucional'' o ``anular'' una norma.

\begin{important}
En inglés, \textit{court} se refiere a cualquier tipo de tribunal, independientemente de su nivel jerárquico. En español, en cambio, se hace una distinción: ``juzgado'' (primera instancia), ``tribunal'' o ``cámara'' (apelación) y ``corte'' (máxima instancia). Cuando aparece \textit{court} en un texto en inglés, debe analizarse el contexto para determinar de qué nivel se trata. Traducir automáticamente \textit{court} como ``corte'' es un error frecuente que puede llevar a confusión.
\end{important}

\subsection{Los \textit{justices} y la distinción con \textit{judge}}
\label{sec:justice}

La siguiente oración del texto trabajado en clase generó confusión en su lectura inicial:

\begin{quote}
\textit{``The justices of the Supreme Court, who can overturn unconstitutional laws, are nominated by the President and confirmed by the Senate.''}
\end{quote}

Una lectura sin atención produce un resultado absurdo: ¿cómo puede ``la justicia'' de la Corte ser nominada por el Presidente? Varias señales alertan: el verbo en plural, las comas que funcionan como inciso aclaratorio, y la lógica jurídica del enunciado.

La resolución: \textit{justice} en plural (\textit{justices}) designa a los \textbf{ministros} de la Corte Suprema, no al concepto abstracto de ``justicia''. El término específico para un juez en general es \textit{judge}. La máxima autoridad de la Corte Suprema recibe el nombre de \textit{Chief Justice}: equivalente funcional al presidente de la Corte Suprema argentina.

% ─────────────────────────────────────────────────────────
\chapter{El sistema judicial federal de los Estados Unidos}
\label{cap:poder-judicial}
% ─────────────────────────────────────────────────────────

\section{La polisemia de \textit{law}}

El término \textit{law} es polisémico: tiene más de un significado según el contexto. Identificar cuál corresponde en cada situación es una tarea central de la lectocomprensión jurídica.

\begin{table}[H]
\centering
\caption{Los distintos significados de \textit{law} según el contexto}
\begin{tabularx}{\textwidth}{>{\raggedright\arraybackslash}p{0.28\textwidth} >{\raggedright\arraybackslash}X >{\raggedright\arraybackslash}p{0.28\textwidth}}
\toprule
\textbf{Significado} & \textbf{Explicación} & \textbf{Ejemplo} \\
\midrule
Una norma en particular & Una ley específica dentro de un sistema jurídico. & \textit{the law of business associations} \\
Un sistema jurídico & El conjunto de normas, decisiones y principios de un ordenamiento. & \textit{English law} \\
Un área del derecho & Una materia específica. & \textit{family law}, \textit{criminal law} \\
La profesión / las fuerzas de seguridad & De forma informal: abogados, jueces, fiscales o incluso la policía. & ``ahí viene la ley'' (la policía) \\
\bottomrule
\end{tabularx}
\end{table}

\section{Falsos cognados del sistema judicial}
\label{sec:court}

\subsection{\textit{Court}: no siempre es ``corte''}

\textit{Court} es el término genérico para tribunal, pero su traducción al español depende del nivel jerárquico:

\begin{itemize}
    \item Un tribunal de primera instancia es un \textbf{juzgado} o \textbf{tribunal}.
    \item Un tribunal de alzada es una \textbf{cámara}.
    \item La instancia máxima es la \textbf{Corte Suprema}.
\end{itemize}

Traducir siempre ``corte'' es un error frecuente, sobre todo cuando se trata de un tribunal de primera instancia.

\begin{example}{Superior Court en Nueva York}
Un tribunal identificado como ``Superior Court'' puede resultar engañoso. Un divorcio de trámite común no llegaría, en el sistema comparado, a una instancia de alzada. Ese ``Superior Court'' era, en realidad, un tribunal de primera instancia. Ante nombres de ese tipo conviene consultar el organigrama judicial oficial de cada estado antes de traducir.
\end{example}

\subsection{\textit{Sentence}: pena, no sentencia}

\begin{definition}{\textit{Sentence} vs. \textit{sentencia}}
En el lenguaje cotidiano, \textit{sentence} significa ``oración'' (gramatical). En el lenguaje jurídico, refiere específicamente a la \textbf{pena} impuesta a quien ya fue declarado culpable: prisión, multa, pena de muerte. Solo existe cuando hubo condena.

La palabra española ``sentencia'' ---que puede ser condenatoria o absolutoria--- se traduce al inglés como \textbf{\textit{judgement}} o \textbf{\textit{resolution}}. La \textit{US Sentencing Commission} es la comisión federal que fija lineamientos para la imposición de penas, no para la determinación de responsabilidad.
\end{definition}

\subsection{\textit{Jurisdiction}: jurisdicción y competencia en un solo término}

En Argentina, \textbf{jurisdicción} es la potestad que poseen todos los jueces por razón de su investidura para decir el derecho. La \textbf{competencia}, en cambio, es la medida de esa jurisdicción: determina a qué juez específico le corresponde entender en cada causa, según la materia, el territorio, el grado o la persona.

En inglés existe una única palabra ---\textit{jurisdiction}--- para ambos conceptos. La distinción práctica para la lectura de textos jurídicos:

\begin{itemize}
    \item En textos de doctrina: \textit{jurisdiction} = ``jurisdicción''.
    \item En la práctica (demandas, fallos, textos procesales): \textit{jurisdiction} = ``competencia''.
\end{itemize}

Cuando la Corte Suprema tiene ``\textit{appellate jurisdiction}'' en la mayoría de las causas, se está hablando de su competencia por apelación. Cuando tiene ``\textit{original jurisdiction}'' en algunos casos excepcionales, se habla de su competencia originaria ---comparable a la competencia originaria de la Corte Suprema argentina en ciertos conflictos entre provincias.

\section{Comparativos y superlativos en el lenguaje judicial}

En la descripción del sistema judicial aparecen con frecuencia comparativos y superlativos que describen la jerarquía de los tribunales.

\begin{table}[H]
\centering
\caption{Comparativos y superlativos en el sistema judicial}
\begin{tabularx}{\textwidth}{>{\raggedright\arraybackslash}p{0.22\textwidth} >{\raggedright\arraybackslash}p{0.18\textwidth} >{\raggedright\arraybackslash}X}
\toprule
\textbf{Término} & \textbf{Tipo} & \textbf{Uso} \\
\midrule
\textbf{highest} & Superlativo & La Corte Suprema: no hay ningún tribunal por encima de ella. \\
\textbf{higher / lower} & Comparativos & Comparan dos instancias entre sí (alzada vs. primera instancia). \\
\textbf{upper / inferior} & Comparativos & Ubican un tribunal respecto de otro; siempre relativos. \\
\bottomrule
\end{tabularx}
\end{table}

``Lower'' y ``upper'' son siempre relativos: un tribunal puede ser ``inferior'' respecto de la Corte Suprema sin ser de primera instancia (por ejemplo, una cámara de apelaciones es ``inferior'' a la Corte Suprema, pero ``superior'' a un juzgado).

\section{La Corte Suprema de Justicia de los Estados Unidos}

La Corte Suprema está integrada por \textbf{nueve miembros}: un \textit{Chief Justice} (presidente de la Corte) y ocho \textit{Associate Justices} (jueces asociados o ministros).

Para resolver una causa se necesita un \textbf{quórum de seis jueces} (\textit{there must be a quorum of six justices}). El uso de \textit{must} señala que se trata de una obligación: si no están presentes los seis jueces, la Corte no puede resolver.

\begin{important}
El cargo de juez de la Corte Suprema es \textbf{vitalicio}. No existe un mandato con plazo fijo: los jueces permanecen en funciones salvo jubilación voluntaria, renuncia, fallecimiento o remoción por las causas establecidas. En promedio, los jueces permanecen dieciséis años en el cargo. Si la Corte trabaja con menos de nueve jueces y se produce un empate, no puede revocar la decisión del tribunal inferior, que queda firme.
\end{important}

\subsection{Designación de los jueces}

El proceso de designación comprende dos pasos: el Presidente propone (nomina) un candidato y el Senado vota para confirmarlo. Habitualmente los candidatos son jueces federales con carrera previa en el sistema.

\section{El sistema de tribunales federales}
\label{sec:tribunales-federales}

La Constitución otorga al Congreso la potestad de crear tribunales federales más allá de la Corte Suprema, para atender causas vinculadas con leyes federales: impositivas, de concursos y quiebras, causas donde sea parte el Estado federal o los gobiernos estaduales, y causas que involucren la Constitución.

\subsection{La distinción entre tribunales federales y estaduales}

El indicador ``\textbf{US}'' (United States) adelante del nombre de un tribunal o institución señala que es federal. Si el nombre indica el nombre de un estado, se trata de un tribunal estadual con su propia organización. Cada estado tiene su propio organigrama judicial, sus propias denominaciones y, en muchos casos, sus propias leyes.

\begin{note}
La palabra \textit{state} es conflictiva para la traducción al español. ``Estado'' puede referirse tanto al gobierno federal como a los gobiernos de cada estado. En Argentina, la distinción equivalente es entre ``estatal'' (federal) y ``provincial''. El grupo de traductores que intentó instalar el término ``estadual'' para los gobiernos de cada estado no logró su objetivo porque el término no existe en español y tampoco resultaba claro.
\end{note}

\subsection{Distintas denominaciones para los mismos tribunales}

En inglés existen distintas formas de nombrar a los tribunales según el criterio de clasificación empleado:

\begin{table}[H]
\centering
\caption{Denominaciones de tribunales federales en inglés}
\begin{tabularx}{\textwidth}{>{\raggedright\arraybackslash}p{0.32\textwidth} >{\raggedright\arraybackslash}p{0.22\textwidth} >{\raggedright\arraybackslash}X}
\toprule
\textbf{Denominación} & \textbf{Instancia} & \textbf{Criterio de denominación} \\
\midrule
\textit{First instance court} & Primera instancia & Nombra directamente la instancia. \\
\textit{District court} & Primera instancia & Nombra la subdivisión territorial (distrito). \\
\textit{Trial court} & Primera instancia & Nombra la función: es donde se produce la prueba. \\
\textit{Court of Appeals / Appellate court} & Alzada & Nombra la función de revisión. \\
\textit{Circuit court} & Alzada & Nombra la subdivisión territorial (circuito). \\
\textit{Upper court / Reviewing court} & Alzada & Nombra la jerarquía o la función de revisar. \\
\bottomrule
\end{tabularx}
\end{table}

El sistema federal también comprende organismos especializados que no son estrictamente tribunales decisorios, como la \textit{Administrative Office of the US Courts} (funciones administrativas), los \textit{Bankruptcy Courts} (concursos y quiebras), la \textit{Court of International Trade}, el \textit{US Tax Court} (materia impositiva) y la \textit{US Sentencing Commission} (lineamientos para la imposición de penas).

\section{El proceso judicial: de los hechos al fallo}

\subsection{La división de funciones entre jurado y juez}

En el sistema anglosajón, los juicios son mayormente por jurado (\textit{trial by jury}). Esta característica introduce una división de funciones que no tiene equivalente en la tradición argentina:

\begin{table}[H]
\centering
\caption{División de funciones entre jurado y juez}
\begin{tabularx}{\textwidth}{>{\raggedright\arraybackslash}p{0.16\textwidth} >{\raggedright\arraybackslash}X >{\raggedright\arraybackslash}X}
\toprule
\textbf{Órgano} & \textbf{Qué resuelve} & \textbf{Cómo se expresa} \\
\midrule
Jurado (\textit{jury}) & Los hechos, a partir de la prueba producida. & En penal: culpable / no culpable (\textit{guilty / not guilty}). \\
Juez (\textit{judge}) & El derecho: cómo se aplica la ley a los hechos. La cuantía de la pena dentro del máximo y el mínimo legal. & Impone la \textit{sentence} (pena) o resuelve la cuestión de fondo en materia civil. \\
\bottomrule
\end{tabularx}
\end{table}

El juez no puede revisar lo resuelto por el jurado en cuanto a los hechos, salvo que mediara un vicio procesal (por ejemplo, un jurado que recibió dinero). También existe el \textit{bench trial}: un juicio en el que no hay jurado y el juez resuelve tanto los hechos como el derecho.

\begin{note}
En el sistema argentino, hasta que se popularizó el juicio por jurados, el juez evaluaba los hechos y el derecho en forma conjunta. El juicio por jurado se encuentra en proceso de implementación gradual en Argentina, lo que hace especialmente relevante la comprensión de este sistema para los futuros profesionales del derecho.
\end{note}

\subsection{El \textit{Discovery} y el \textit{hearsay}}

En el sistema estadounidense existe una instancia previa al juicio con jurado denominada \textit{Discovery}. No tiene equivalente directo en el sistema argentino.

\begin{definition}{\textit{Discovery}}
El \textit{Discovery} es una etapa previa al juicio en la que se produce prueba anticipada: los testigos declaran y se anticipa la evidencia. El juez no está presente físicamente de manera constante, pero las partes pueden acudir a él ante cualquier anomalía. Las partes van preparando su estrategia: si un testimonio puede perjudicar su posición, deciden cómo abordarlo antes de llegar al jurado.
\end{definition}

\begin{definition}{\textit{Hearsay}}
Un testigo debe declarar sobre hechos de los que tuvo conocimiento directo. El testimonio basado en lo que un tercero contó (\textit{``me dijo la vecina, que le dijo el primo''}) se denomina \textit{hearsay} y no es admisible como prueba. Si una parte intenta presentar un testigo con estas características, la contraparte puede recurrir al juez para impugnarlo durante el \textit{Discovery}.
\end{definition}

\subsection{Sinónimos de ``juicio'' en inglés}

El inglés dispone de varios términos que se traducen aproximadamente como ``juicio'', ``causa'' o ``proceso'', sin ser sinónimos exactos entre sí:

\begin{table}[H]
\centering
\caption{Sinónimos de ``juicio'' en inglés jurídico}
\begin{tabularx}{\textwidth}{>{\raggedright\arraybackslash}p{0.22\textwidth} >{\raggedright\arraybackslash}X}
\toprule
\textbf{Término} & \textbf{Uso y observación} \\
\midrule
\textit{lawsuit} & El más general; de uso bastante informal. Deriva de \textit{sue} (demandar). \\
\textit{trial} & Proceso en sentido amplio (\textit{to go to trial}); en sentido técnico, la etapa probatoria. \\
\textit{case} & Causa; caso. \textit{Case number}: número de causa (expediente). \\
\textit{action / claim} & El acto que da inicio al proceso (\textit{to bring an action}). \\
\textit{proceeding(s)} & Proceso; abarca las distintas actuaciones dentro de la causa. \\
\bottomrule
\end{tabularx}
\end{table}

\section{El camino de apelación hacia la Corte Suprema}

La mayoría de las causas llegan a la Corte Suprema por vía de apelación (\textit{appellate jurisdiction}). Solo excepcionalmente tiene competencia originaria (\textit{original jurisdiction}), como en conflictos entre estados, comparables a conflictos entre provincias en el sistema argentino.

Para que la Corte acepte revisar una causa, se necesitan \textbf{cuatro votos de los nueve jueces} (\textit{four justices must vote in favor for a case to be granted review}). La Corte recibe entre 7000 y 8000 solicitudes de revisión por año y concede entre 70 y 80, lo que refleja la extrema selectividad del tribunal.

Una vez concedida la revisión, las partes presentan \textbf{\textit{briefs}} (escritos con fundamentos y argumentos) y realizan argumentos orales. Cada parte dispone de treinta minutos para exponer su caso.

\subsection{Votos de mayoría y disidencia}

Tras los argumentos orales, cada juez elabora su voto (\textit{opinion}). La mayoría de los votos determina la decisión de la Corte:

\begin{table}[H]
\centering
\caption{Tipos de votos en la Corte Suprema}
\begin{tabularx}{\textwidth}{>{\raggedright\arraybackslash}p{0.30\textwidth} >{\raggedright\arraybackslash}p{0.28\textwidth} >{\raggedright\arraybackslash}X}
\toprule
\textbf{Concepto en inglés} & \textbf{Traducción} & \textbf{Explicación} \\
\midrule
\textit{Opinion} & Voto & Fundamentación individual de cada juez. \\
\textit{Majority opinion} & Voto de la mayoría & Más de la mitad de los jueces en el mismo sentido: se convierte en la decisión de la Corte. \\
\textit{Dissenting / minority opinion} & Voto en disidencia & Voto de quienes no acompañan la posición mayoritaria. \\
\textit{Hands down a decision} & Dictar sentencia & La resolución, una vez emitida, ya no puede modificarse. \\
\bottomrule
\end{tabularx}
\end{table}

\section{Comparación final: Estados Unidos y Argentina}

\begin{table}[H]
\centering
\caption{Diferencias estructurales entre los sistemas judiciales}
\begin{tabularx}{\textwidth}{>{\raggedright\arraybackslash}p{0.26\textwidth} >{\raggedright\arraybackslash}X >{\raggedright\arraybackslash}X}
\toprule
\textbf{Aspecto} & \textbf{Estados Unidos} & \textbf{Argentina} \\
\midrule
Leyes de fondo & Cada estado tiene sus propias leyes. Un mismo hecho puede ser delito en un estado y no en otro. & Códigos de fondo únicos para todo el país. Las leyes procesales varían por provincia. \\
Pena de muerte & Existe en algunos estados y en otros no. & No está prevista en el ordenamiento vigente. \\
Juicio por jurado & Es la regla. El jurado resuelve los hechos; el juez, el derecho. & Implementación gradual; históricamente el juez resolvía tanto hechos como derecho. \\
\bottomrule
\end{tabularx}
\end{table}

% ─────────────────────────────────────────────────────────
\chapter{La Constitución del Reino Unido}
\label{cap:reino-unido}
% ─────────────────────────────────────────────────────────

\section{Contexto geográfico y político: Reino Unido e Inglaterra}

Antes de avanzar en el análisis constitucional, conviene precisar una confusión frecuente de cultura general.

\begin{important}
El nombre oficial del país es \textit{United Kingdom of Great Britain and Northern Ireland}. \textit{Great Britain} abarca Inglaterra, Gales y Escocia. Inglaterra es, por lo tanto, solo una parte del Reino Unido. Irlanda del Norte forma parte del Reino Unido, mientras que Irlanda (del Sur) queda por fuera.

El gentilicio para el Reino Unido en su conjunto es \textit{British}. Este fenómeno es comparable al uso de ``americano'' para referirse exclusivamente a los ciudadanos de los Estados Unidos: una apropiación lingüística del término colectivo por parte del país dominante.
\end{important}

\section{¿Qué es una Constitución? Las tres funciones generales}

El texto trabajado en clase comienza con una definición general de Constitución aplicable a cualquier sistema:

\begin{formula}{Las tres funciones de toda Constitución}
\textit{``Constitutions organize, distribute and regulate state power.''}

Las constituciones \textbf{organizan} el poder estatal (establecen las principales instituciones del Estado), lo \textbf{distribuyen} (entre los distintos poderes o territorios) y \textbf{regulan} las relaciones entre esas instituciones y entre ellas y la ciudadanía.
\end{formula}

\section{\textit{Unwritten} vs. \textit{uncodified}: una distinción clave}

El texto utiliza el término \textit{unwritten} entre comillas, lo que es una señal para el lector: el término resulta impreciso. El Reino Unido sí cuenta con normas escritas; lo que no tiene es un único instrumento donde estén codificadas.

\begin{important}
\textbf{No escrita $\neq$ no codificada.} El Reino Unido sí tiene normas escritas: las leyes del Parlamento son textos escritos. Lo que no existe es un único documento codificado donde esté reunido el contenido constitucional, a diferencia de sistemas como el argentino (Constitución Nacional de 1853) o el estadounidense (Constitución de 1787).
\end{important}

Solo tres países modernos carecen de constitución codificada en un único documento: Reino Unido, Nueva Zelanda e Israel.

\begin{table}[H]
\centering
\caption{Ventajas y desventajas de una constitución no codificada}
\begin{tabularx}{\textwidth}{>{\raggedright\arraybackslash}X >{\raggedright\arraybackslash}X}
\toprule
\textbf{Ventaja} & \textbf{Desventaja} \\
\midrule
Mayor facilidad para modificar el sistema, precisamente por no estar concentrado en un único instrumento que exija un procedimiento agravado. & Mayor dificultad de acceso y sistematización para quienes deben invocar o fundamentar un principio constitucional. \\
\bottomrule
\end{tabularx}
\end{table}

\section{Evolución frente a revolución: la estabilidad histórica británica}

A diferencia de países que debieron ``empezar de cero'' (\textit{start from scratch}) tras un cambio de régimen ---como Argentina tras su independencia, que debió construir instituciones prácticamente desde cero---, el sistema constitucional británico evolucionó gradualmente a lo largo de un período extenso, sin una ruptura institucional equivalente.

Esta ausencia de ruptura explica la estabilidad relativa del sistema político británico (\textit{the relative stability of the British polity}), que a su vez tiene consecuencias concretas para el derecho comercial internacional.

\begin{example}{El derecho inglés en el comercio internacional}
El derecho mercantil vinculado a la compraventa internacional de mercaderías se rige, en gran medida, por el \textit{Sale of Goods Act}. La regulación internacional de contratos elige con frecuencia el derecho inglés porque se ha mantenido sustancialmente estable desde el siglo XV. En seguros marítimos, la tradición de \textit{Lloyd's} ---la primera compañía de seguros marítimos--- constituye un antecedente que el comercio internacional, incluido el estadounidense, terminó adoptando. Esta estabilidad hace al derecho inglés especialmente atractivo para el comercio internacional, que requiere previsibilidad.
\end{example}

\section{Falsos cognados jurídicos fundamentales del sistema británico}
\label{sec:falsos-cognados-uk}

Tres falsos cognados son particularmente peligrosos en la lectura de textos sobre el sistema constitucional del Reino Unido. Un error en cualquiera de ellos lleva a una comprensión distorsionada de todo el texto.

\subsection{\textit{Jurisprudence}: filosofía del derecho, no fallos judiciales}

\begin{important}
\textit{Jurisprudence} en inglés no significa ``jurisprudencia'' en el sentido de fallos judiciales. Refiere a la \textbf{filosofía del derecho} o a la \textbf{doctrina}, en un sentido próximo a la teoría del derecho. Para referirse a los fallos judiciales se emplean: \textit{legal precedent}, \textit{case law} o \textit{judicial decisions}.
\end{important}

Esta distinción tiene consecuencias directas para la comprensión de la jerarquía de fuentes: en el texto sobre la Constitución del Reino Unido, cuando se menciona \textit{jurisprudence} no se está hablando de los precedentes judiciales (\textit{case law}), sino de la doctrina académica.

\subsection{\textit{Statute}: ley del Parlamento, no estatuto}

\begin{important}
En el derecho inglés, \textit{statute} no equivale a ``estatuto'' en el sentido de norma interna de una institución. Un \textit{statute} es una \textbf{ley formal y escrita, sancionada por el Parlamento}, equivalente a lo que en español llamamos simplemente ``ley''. No debe confundirse con el estatuto universitario ni con el estatuto social de una sociedad comercial.
\end{important}

Los \textit{statutes} constituyen la fuente jerárquicamente más alta dentro del sistema constitucional británico.

\subsection{\textit{Convention}: costumbre institucional, no convención internacional}

\begin{important}
Una \textit{convention}, en el sistema constitucional británico, no es un tratado ni una convención internacional. Se trata de \textbf{prácticas institucionales no escritas} que regulan el ejercicio del gobierno (\textit{the business of governing}), consolidadas por la tradición a lo largo del tiempo. Es, en definitiva, costumbre institucional.
\end{important}

\begin{example}{\textit{Convention} en la práctica}
El Primer Ministro se reúne periódicamente con la Corona para sesiones informativas (\textit{briefings}). Esta práctica no consta en ninguna ley escrita; es una \textit{convention}: una costumbre que ``lubrica'' el engranaje del sistema constitucional.
\end{example}

El siguiente cuadro consolida los tres falsos cognados del sistema constitucional británico:

\begin{table}[H]
\centering
\caption{Falsos cognados jurídicos del sistema constitucional británico}
\begin{tabularx}{\textwidth}{>{\raggedright\arraybackslash}p{0.22\textwidth} >{\raggedright\arraybackslash}p{0.30\textwidth} >{\raggedright\arraybackslash}X}
\toprule
\textbf{Término en inglés} & \textbf{Falso sentido} & \textbf{Significado correcto} \\
\midrule
\textit{jurisprudence} & Jurisprudencia (fallos) & Filosofía del derecho; doctrina \\
\textit{statute} & Estatuto (norma interna) & Ley formal del Parlamento \\
\textit{convention} & Convención internacional & Costumbre institucional no escrita \\
\bottomrule
\end{tabularx}
\end{table}

\section{Las fuentes de la Constitución británica}

En lugar de un único documento codificado, la Constitución del Reino Unido se nutre de varias fuentes que coexisten. Para visualizar cómo cambia la jerarquía de fuentes respecto del sistema continental, puede adaptarse didácticamente la pirámide de Kelsen ---con la aclaración expresa de que se trata de un recurso pedagógico y no de una construcción reconocida por la doctrina constitucionalista:

\begin{table}[H]
\centering
\caption{Comparación de la jerarquía de fuentes: sistema continental vs. sistema anglosajón}
\begin{tabularx}{\textwidth}{>{\raggedright\arraybackslash}p{0.22\textwidth} >{\raggedright\arraybackslash}p{0.36\textwidth} >{\raggedright\arraybackslash}X}
\toprule
\textbf{Nivel} & \textbf{Sistema continental} & \textbf{Sistema anglosajón (Reino Unido)} \\
\midrule
Superior & Constitución y tratados con jerarquía constitucional & \textit{Statutes}: leyes del Parlamento \\
Intermedio & Leyes ordinarias & \textit{Conventions}: costumbre institucional \\
Relevante & Jurisprudencia y costumbre (escasa relevancia formal) & \textit{Common law / case law}: altísima relevancia como fuente \\
Consultivo & --- & \textit{Constitutional authorities}: doctrina y especialistas \\
\bottomrule
\end{tabularx}
\end{table}

\subsection{Los \textit{statutes}: leyes del Parlamento}

Los \textit{statutes} son las normas que emanan del Parlamento: \textit{``laws passed by Parliament [...] the highest form of law''}. Constituyen la fuente jerárquicamente más alta.

\subsection{Las \textit{conventions}: costumbre institucional}

Las \textit{conventions} son prácticas no escritas que regulan el ejercicio del gobierno, consolidadas por tradición. Determinadas potestades quedan reservadas por costumbre a ciertas instituciones sin que exista una norma escrita que así lo establezca.

\subsection{El \textit{common law}: decisiones judiciales}

\begin{example}{Explicación del \textit{common law}}
El \textit{common law} es \textit{``law developed by the courts and judges through cases''}: el derecho desarrollado por los tribunales a través de los casos. Se aplica por analogía: la solución adoptada en un caso se extiende a causas con identidad de circunstancias. Ocupa un lugar central en el sistema anglosajón, a diferencia del sistema continental donde la jurisprudencia tiene un peso formal mucho menor.
\end{example}

\subsection{El derecho de la Unión Europea y el Brexit}

El texto de trabajo sobre el sistema constitucional británico resulta parcialmente desactualizado en este punto: hace referencia a la incorporación de normas europeas al ordenamiento británico, proceso superado tras la salida del Reino Unido de la Unión Europea. El Brexit implicó que el Reino Unido debió analizar qué normas europeas conservar dentro de su ordenamiento interno, sin que ello quedara plasmado en un instrumento constitucional único. El propio Brexit constituye un ejemplo de la flexibilidad del sistema.

\subsection{El derecho internacional y las \textit{constitutional authorities}}

El texto agrega que el Reino Unido está sujeto al derecho internacional (\textit{UK is also subject to international law}). Ante la ausencia de un texto único al que recurrir, ciertas autoridades y especialistas ---jueces, abogados y doctrinarios--- cumplen una función consultiva como \textit{constitutional authorities}. Se trata de una función asimilable a una recomendación, razón por la cual se ubica en el nivel más bajo de la pirámide comparada.

\section{Soberanía parlamentaria y principios centrales del sistema}

\begin{formula}{El principio organizador del sistema británico}
\textit{``What the Queen in Parliament enacts is law.''}

Traducción: lo que la Corona en el Parlamento promulga es ley. El texto usa \textit{``the Queen''} por corresponder a una versión anterior al cambio de titularidad de la Corona.
\end{formula}

A diferencia del sistema argentino, donde la potestad legislativa corresponde al Congreso porque se la otorga la Constitución, en el Reino Unido es la Corona quien históricamente delega esa potestad al Parlamento. Los legisladores son, literalmente, ``hacedores de leyes'' (\textit{law-makers}).

\begin{definition}{Soberanía parlamentaria (\textit{parliamentary sovereignty})}
El principio definitorio de la Constitución británica según el cual el Parlamento, democráticamente electo, posee el poder legislativo último (\textit{ultimate law-making power}) para crear o abolir cualquier ley. El carácter ``último'' no remite a un orden jerárquico dentro de un ordenamiento, sino al carácter definitivo de esa potestad, que recae en el centro del sistema unitario: el Parlamento.
\end{definition}

Los otros principios centrales identificados en el texto son: el \textit{rule of law} (estado de derecho o ``imperio de la ley''; no debe traducirse literalmente como ``regla de la ley''), la separación entre los poderes ejecutivo, legislativo y judicial, y la existencia de un Estado unitario en el cual el poder último recae en el Parlamento.

\subsection{El carácter ``mítico'' de algunos principios}

El texto introduce mediante el conector \textit{however} (sin embargo) un matiz crítico:

\begin{important}
Pese a proclamarse la separación de poderes, en la práctica existe una fusión entre el ejecutivo y el legislativo, sostenida por la tradición histórica más que por una estructura de Estado moderno estrictamente diferenciada. El término \textit{mythical} no remite a algo inexistente, sino a una construcción sostenida por la tradición que no encaja del todo con la estructura de un Estado moderno.

La soberanía parlamentaria como principio absoluto fue además cuestionada por la incorporación de normativa europea previa al Brexit: el Parlamento debió incorporar legislación de instituciones europeas ajenas a su propia estructura soberana.
\end{important}

\section{Hitos históricos del constitucionalismo británico}

\subsection{La Magna Carta (1215)}

La Magna Carta (\textit{The Great Charter of the Liberties of England}) constituye la primera declaración de libertades en la historia constitucional moderna. Estableció que los gobernantes no podían hacer lo que quisieran arbitrariamente y debían estar sujetos a la ley. Su importancia radica no tanto en sus disposiciones específicas ---muchas fueron modificadas--- como en la consagración del principio de limitación del poder mediante ley. En el vocabulario del texto: constituye una \textbf{\textit{landmark law}}, un hito legal que marcó un punto de inflexión.

\begin{definition}{Magna Carta (1215)}
Primera declaración de libertades de Inglaterra, que estableció que incluso los gobernantes están sujetos a la ley. Sentó las bases conceptuales para el gobierno constitucional y la libertad bajo el imperio de la ley.
\end{definition}

\subsection{El Bill of Rights (1689)}

El Bill of Rights de 1689 fue sancionado tras la Revolución Gloriosa y consolidó jurídicamente la supremacía parlamentaria sobre las prerrogativas monárquicas. Estableció: la supremacía del Parlamento sobre la Corona; sesiones regulares del Parlamento; elecciones libres para la Cámara Baja; libertad de expresión en debates parlamentarios; y prohibición de castigos crueles e inusuales.

\begin{important}
Es fundamental distinguir el Bill of Rights británico de 1689 del estadounidense de 1791. El documento británico versaba sobre derechos políticos del Parlamento frente a la Corona; el estadounidense se enfocaba en derechos fundamentales de los ciudadanos. Ambos son documentos constitucionales fundamentales, pero con objetos de protección completamente distintos.
\end{important}

\section{Problemas y reformas constitucionales}

El texto identifica dos problemas derivados de la falta de codificación:

\begin{itemize}
    \item \textbf{Falta de certeza:} al encontrarse la Constitución dispersa en múltiples fuentes, resulta difícil identificar dónde reside el contenido constitucional aplicable a un caso.
    \item \textbf{Mayor facilidad de reforma:} al no existir un instrumento codificado con estatus de ley superior que exija un procedimiento agravado, la modificación del sistema es relativamente sencilla.
\end{itemize}

Desde 1997 se incorporaron reformas sucesivas: la abolición de la mayoría de las bancas hereditarias en la Cámara de los Lores (\textit{House of Lords}), la incorporación del \textit{Human Rights Act} y los \textit{devolution settlements} con Escocia, Gales e Irlanda del Norte.

\begin{definition}{\textit{Devolution settlements}}
Acuerdos mediante los cuales el gobierno central delega potestades propias a los gobiernos de Escocia, Gales e Irlanda del Norte. A diferencia de lo que ocurre en un sistema federal, esta delegación es \textbf{revocable}: podría en teoría ser derogada por el propio Reino Unido. Esto genera tensiones recurrentes respecto de la extensión y la reversibilidad de las potestades delegadas.
\end{definition}

% ─────────────────────────────────────────────────────────
% CUADRO CORNELL — PARTE II
% ─────────────────────────────────────────────────────────
\chapter*{Cuadro Cornell: Derecho constitucional comparado}
\addcontentsline{toc}{chapter}{Cuadro Cornell: Derecho constitucional comparado}
\markboth{Cuadro Cornell: Derecho constitucional comparado}{Cuadro Cornell: Derecho constitucional comparado}

\begin{longtable}{
    >{\raggedright\arraybackslash}p{0.30\textwidth}
    >{\raggedright\arraybackslash}p{0.60\textwidth}
}
\toprule
\textbf{Preguntas / palabras clave} & \textbf{Notas / respuestas} \\
\midrule
\endfirsthead
\toprule
\textbf{Preguntas / palabras clave} & \textbf{Notas / respuestas} \\
\midrule
\endhead
\textbf{¿Qué es \textit{government} en el texto constitucional?} &
El Estado en su conjunto (ejecutivo + legislativo + judicial), no el poder ejecutivo. Es un falso cognado respecto del uso argentino de ``gobierno'' (que designa solo al ejecutivo). \\

\textbf{¿Cuántos senadores y representantes tiene el Congreso?} &
Senado: 100 senadores (2 por estado, independientemente de la población); mandato de 6 años; reelección ilimitada. Cámara de Representantes: 435 miembros, distribuidos proporcionalmente a la población; mandato de 2 años; reelección ilimitada. \\

\textbf{¿Qué significa \textit{term} en un contexto político?} &
Mandato; período de ejercicio de un cargo. No ``término'' en sentido de palabra. Ejemplo: ``a four-year term'' = un mandato de cuatro años. \\

\textbf{¿Cuál es el mandato presidencial en EE.UU. y cuántas veces puede reelegirse?} &
Mandato de 4 años; máximo dos mandatos (8 años). Límite establecido constitucionalmente a partir de la presidencia de Franklin D. Roosevelt. \\

\textbf{¿Qué es la Corte Suprema y cómo se designan sus jueces?} &
9 miembros: 1 Chief Justice + 8 Associate Justices. Designación: nominación presidencial + confirmación del Senado. Cargo vitalicio (salvo jubilación, renuncia, fallecimiento o remoción). Quórum: 6 jueces. \\

\textbf{¿Cómo llega un caso a la Corte Suprema?} &
Mayormente por apelación (necesita 4 votos de 9). Rara vez por competencia originaria. La Corte recibe entre 7000 y 8000 solicitudes por año y concede entre 70 y 80. \\

\textbf{¿Qué diferencia hay entre \textit{court}, ``corte'', ``cámara'' y ``juzgado''?} &
\textit{Court} es el término genérico en inglés para cualquier tribunal. En español: juzgado (primera instancia), cámara (alzada), corte (máxima instancia). Traducir siempre ``corte'' es un error frecuente. \\

\textbf{¿Qué es el \textit{Discovery}?} &
Etapa previa al juicio con jurado en EE.UU., sin equivalente argentino. Se produce prueba anticipada y las partes preparan su estrategia. El \textit{hearsay} (testimonio de oídas) no es admisible como prueba. \\

\textbf{¿Qué diferencia hay entre \textit{unwritten} y \textit{uncodified}?} &
El RU sí tiene normas escritas (statutes del Parlamento). Lo que no tiene es un único instrumento codificado. ``Unwritten'' (no escrita) es impreciso; ``uncodified'' (no codificada) es lo correcto. \\

\textbf{¿Cuáles son las tres funciones de toda Constitución?} &
Organizar, distribuir y regular el poder estatal (organize, distribute and regulate state power). \\

\textbf{¿Qué es \textit{statute} en el derecho inglés?} &
Ley formal y escrita sancionada por el Parlamento. No debe confundirse con ``estatuto'' (norma interna de una institución u organización). \\

\textbf{¿Qué son las \textit{conventions}?} &
Prácticas institucionales no escritas que regulan el ejercicio del gobierno, consolidadas por la costumbre. No son convenciones internacionales. \\

\textbf{¿Qué es la soberanía parlamentaria?} &
El principio definitorio de la Constitución británica: el Parlamento, democráticamente electo, posee el poder legislativo último para crear o abolir cualquier ley. \\

\textbf{¿Qué son los \textit{devolution settlements}?} &
Acuerdos por los cuales el gobierno central delega potestades a Escocia, Gales e Irlanda del Norte. A diferencia de un sistema federal, esta delegación es revocable. \\

\textbf{¿Qué estableció la Magna Carta (1215) y el Bill of Rights (1689)?} &
Magna Carta: los gobernantes están sujetos a la ley; bases del gobierno constitucional. Bill of Rights (1689): supremacía parlamentaria sobre la Corona; garantías de sesiones regulares, elecciones libres y prohibición de castigos crueles. Distinto del Bill of Rights estadounidense (1791), que protegía derechos ciudadanos. \\

\midrule
\multicolumn{2}{p{0.90\textwidth}}{
\textbf{Resumen:} La Unidad 1 abordó el derecho constitucional comparado a partir de dos sistemas del \textit{common law}: Estados Unidos y el Reino Unido. Del análisis de la organización estatal estadounidense se extrajeron los primeros falsos cognados (\textit{government}, \textit{justice}), la lógica del sistema de frenos y contrapesos, y la estructura del Congreso, el ejecutivo y el poder judicial federal, incluyendo el proceso de apelación hacia la Corte Suprema y las diferencias del proceso judicial anglosajón (jurado, \textit{Discovery}, \textit{hearsay}). El sistema constitucional del Reino Unido, en cambio, demandó comprender una lógica radicalmente distinta: la de una constitución dispersa en fuentes no codificadas (\textit{statutes}, \textit{conventions}, \textit{common law}) cuya coherencia surge de la tradición histórica más que de un texto único. Los tres falsos cognados (\textit{jurisprudence}, \textit{statute}, \textit{convention}) y la distinción entre lo ``no escrito'' y lo ``no codificado'' son los pilares conceptuales de esa unidad.
} \\
\bottomrule
\end{longtable}

% =====================================================
% PARTE III
% =====================================================
\part{Unidad 2: Derecho contractual comparado}

% ─────────────────────────────────────────────────────────
\chapter{La formación del contrato en el \textit{common law}}
\label{cap:contratos}
% ─────────────────────────────────────────────────────────

El tránsito de la Unidad 1 a la Unidad 2 implica un desplazamiento del derecho público al derecho privado: del Estado y su organización constitucional a la relación entre particulares y la formación de sus acuerdos. La Unidad 2 resulta especialmente relevante para la práctica profesional porque gran parte de los contratos que regulan la vida cotidiana ---\textit{streaming}, aplicaciones, comercio electrónico, alquileres temporarios, seguros de viaje--- están redactados originalmente en inglés bajo las categorías que aquí se estudian.

\begin{important}
Comprender cómo el \textit{common law} construye la formación del contrato permite identificar qué cláusula cumple cada función, anticipar cómo un tribunal extranjero interpretaría un documento y evitar falsos amigos centrales como \textit{consideration} (que no es ``consideración'' sino contraprestación) o \textit{evidence} (que no es ``evidencia'' sino medios de prueba).
\end{important}

\section{Los cuatro elementos de la formación del contrato}

El texto trabajado en clase identifica cuatro elementos que, según el \textit{common law}, integran la formación de un contrato. Puede pensarse la relación entre ellos con un ejemplo cotidiano: alguien contrata a un pintor. Primero debe haber intención real de contratar; luego una oferta concreta y una aceptación en esos mismos términos; después una contraprestación; y finalmente los términos deben ser suficientemente ciertos. Los cuatro elementos son eslabones de una misma cadena: sin cualquiera de ellos, no hay contrato exigible.

\section{La intención de contratar (\textit{intention})}

El primer elemento es la intención de las partes de efectivamente celebrar un contrato vinculante (\textit{only binding relationship}). Una mera propuesta o una idea vaga no constituye un contrato. Ante la duda, el tribunal tomará en cuenta cualquier prueba disponible (\textit{any evidence}) y la naturaleza del vínculo entre las partes (\textit{the nature of the relationship between the parties}).

\begin{definition}{\textit{Means of evidence}}
\textit{Means of evidence} equivale a los ``medios de prueba'' del derecho argentino, no simplemente a ``evidencia''. Este es uno de los falsos amigos más frecuentes del inglés jurídico: \textit{evidence} no es ``evidencia'' sino prueba.
\end{definition}

\subsection{La presunción en vínculos familiares y afectivos}

Como regla general, en las relaciones familiares o entre personas cercanas se presume que \textbf{no} existe intención de celebrar un contrato vinculante. La presunción alcanza a los cónyuges, a vínculos de amistad y a los miembros de una misma familia.

Sin embargo, esta presunción puede ser \textbf{refutada} ---no anulada, sino refutada--- cuando una de las partes demuestra con pruebas que sí existió intención de contratar. El tribunal puede considerar: si existió un acuerdo escrito; si intervinieron abogados; si alguna parte actuó confiando en el acuerdo; y la naturaleza real del vínculo, más allá de la forma o el título que se le haya dado.

\begin{example}{Caso \textit{Merritt v. Merritt} (1970)}
El matrimonio Merritt era propietario de una vivienda hipotecada. Al separarse, firmaron un acuerdo escrito por el cual, si la mujer pagaba la hipoteca íntegramente, el señor Merritt le transferiría la propiedad. Cuando la señora Merritt terminó de pagar, el señor Merritt sostuvo que no había tenido intención de contratar y se negó a transferir. El tribunal resolvió que efectivamente había existido un contrato y ordenó la transferencia de dominio: pese al vínculo conyugal, la señora Merritt pudo aportar prueba de que las partes sí habían tenido intención de contratar.
\end{example}

\subsection{La presunción en vínculos comerciales}

En los vínculos comerciales, la presunción es exactamente inversa: se presume que las partes \textbf{sí} tienen intención de celebrar un contrato legalmente vinculante. Las negociaciones previas (\textit{bargains}) ---correos electrónicos, presupuestos, señas--- se toman como indicios de esa intención.

El término \textit{bargains} merece atención: en registro coloquial se acerca a ``regatear'', pero en sentido jurídico refiere a las negociaciones precontractuales y las tratativas previas. Esta presunción comercial también puede refutarse: demostrando que, pese al vínculo comercial, las partes dejaron expresamente asentado que aquello no se tendría por un contrato.

\begin{note}
Una diferencia relevante con el derecho argentino: en el texto trabajado, el primer elemento de la formación del contrato se enfoca exclusivamente en el tipo de vínculo entre las partes y en cómo se presume o no la intención. La capacidad ---vinculada a la edad, la capacidad mental y el discernimiento--- es analizada en el derecho argentino como un elemento autónomo del consentimiento, no incluido aquí.
\end{note}

\section{Oferta y aceptación (\textit{Offer and Acceptance})}

El segundo elemento es el entendimiento mutuo sobre los términos y condiciones del contrato ---denominado en el texto \textit{meeting of the minds} (acuerdo de voluntades). Para que ese entendimiento exista, deben concurrir dos elementos distintos: una oferta y una aceptación.

\begin{definition}{\textit{Enforceable}}
Exigible o ejecutable: si una parte no cumple, la otra puede reclamar ante un tribunal, pero solo si están dadas las condiciones para ello. Los \textit{terms and conditions} son, en la tradición jurídica argentina, las cláusulas del contrato: el plazo, el precio, la forma de cumplimiento.
\end{definition}

\subsection{Requisitos de la oferta}

El texto define la oferta como \textit{purposeful}: con propósito, con intencionalidad, no una manifestación casual. La oferta debe indicar claramente los términos del contrato, porque los tribunales deben evaluar si esos términos fueron suficientemente precisos.

\begin{definition}{\textit{Offer}}
\textit{``...a proposal, a promise, or other manifestation of willingness to enter into a contract or bargain under proposed terms with another party that has the power to accept it.''}

El sufijo \textit{-ness} sustantiva y vuelve abstracta la palabra: \textit{willingness} = voluntad. \textit{Will} también significa voluntad y, en materia sucesoria, testamento (\textit{last will and testament}).
\end{definition}

\begin{example}{Oferta demasiado vaga}
Si alguien propone celebrar un contrato pero no expresa el monto, no aclara el plazo ni especifica concretamente qué es lo que quiere que se haga, esa propuesta resulta demasiado vaga para constituir una oferta en sentido estricto.
\end{example}

La forma de comunicación de una oferta puede ser variada: lo determinante es cuán claramente estén expresados todos sus términos.

\subsection{La vigencia temporal de las ofertas}

Las ofertas pueden caducar tras un período de tiempo razonable, aunque no se haya fijado un plazo expreso (\textit{offers can expire after a reasonable time period}).

\begin{example}{El collar reclamado 50 años después}
Si alguien ofrece una recompensa por encontrar un collar perdido y una persona se presenta cincuenta años después a reclamarla, sería ilógico considerar que la oferta seguía vigente: la oferta ya no estaba disponible para ser aceptada (\textit{off the table}).
\end{example}

\subsection{La contraoferta (\textit{counter offer})}

La aceptación debe darse en los mismos términos que la oferta. Si la aceptación modifica algún término, no hay aceptación sino una \textbf{contraoferta} (\textit{counter offer}), que deja sin efecto la oferta original y da lugar a una nueva propuesta que la otra parte deberá aceptar.

Al plantearse una contraoferta, los roles se invierten: quien era \textit{offeror} pasa a ser \textit{offeree} de la nueva propuesta, y viceversa. El proceso puede repetirse hasta que exista una aceptación lisa y llana de una oferta concreta.

La aceptación debe comunicarse de manera expresa (oral o escrita) o implícita, a través de conductas o signos inequívocos que permitan entender que se presta el consentimiento.

\section{La contraprestación (\textit{consideration})}

\begin{important}
\textit{Consideration} no significa ``consideración'': es el falso amigo más frecuente del inglés contractual. \textit{Consideration} es la \textbf{contraprestación}: el intercambio de promesas o prestaciones recíprocas entre las partes (\textit{an exchange of promises in which each party promises something to the other}).
\end{important}

\begin{definition}{\textit{Consideration}}
\textit{What each person promises the other is consideration}: aquello que se da a cambio. La contraprestación implica un beneficio para quien la recibe y, correlativamente, un \textit{detriment} ---pérdida o desventaja en sentido amplio--- para quien la entrega.
\end{definition}

\subsection{Algo de valor (\textit{something of value})}

La contraprestación debe involucrar algo de valor: algo susceptible de apreciación pecuniaria, aunque también pueda tener un valor sentimental o simbólico. Un caramelo, por ejemplo, no sería representativo de un intercambio con entidad suficiente para sostener la celebración de un contrato.

Lo que \textit{induces} ---motiva--- a las partes a contratar es justamente el deseo de obtener aquello que la otra parte ofrece a cambio. La fórmula que también utiliza el derecho argentino: nadie celebra un contrato por nada.

\subsection{Formas de la contraprestación}

La contraprestación no siempre consiste en dinero: puede tratarse de un intercambio de favores, de bienes, de servicios, o incluso de la \textbf{abstención} de realizar una determinada conducta (\textit{to refrain from doing something}).

\begin{example}{La parcela del cementerio familiar}
El texto ilustra la contraprestación por abstención con un caso en el que alguien se comprometió a no ocupar una determinada parcela de un cementerio familiar. El tribunal entendió que ese compromiso de no hacer ---sin que mediara pago de dinero--- constituía igualmente una contraprestación válida. El principio tiene vigencia actual: en los contratos online, aceptar términos de servicio sin pagar puede equivaler a ceder derechos (a datos de navegación, publicidad) como contraprestación.
\end{example}

\section{Certeza de los términos (\textit{Certainty of Terms})}

El cuarto elemento exige que los términos del contrato sean claros y completos: ambas partes deben comprender exactamente qué están contratando. Si faltan condiciones, cualquiera puede luego cuestionar ante un tribunal si efectivamente eso fue lo pactado.

\begin{formula}{Ambigüedad en los términos esenciales}
\textit{Contracts may be void/unenforceable if key information in the terms is ambiguous}: si hay ambigüedad o vaguedad en la información esencial del contrato, este puede considerarse inválido o inejecutivo. Los tribunales tenderán a favorecer la validez del contrato e intentarán interpretarlo de manera que sobreviva, pero no redactarán ni completarán lo que las partes omitieron.
\end{formula}

\subsection{Nulidad total o parcial}

\begin{formula}{El límite de la interpretación judicial}
\textit{Courts will not create entire contracts for the parties}: si el tribunal no cuenta con información suficiente para interpretar el contrato, no lo redactará ni completará por las partes. En tal caso, el contrato queda sin efecto, total o parcialmente.
\end{formula}

\begin{definition}{\textit{Unenforceable} vs. \textit{invalid}}
\textit{Unenforceable}: la cláusula o estipulación imprecisa queda sin efecto (nulidad parcial); el resto del contrato puede subsistir. \textit{Invalid}: el contrato en su totalidad resulta nulo.
\end{definition}

\subsection{Interpretación judicial: los usos, las tratativas y ``\textit{the four corners}''}

Antes de declarar la nulidad, el tribunal buscará interpretar el contrato a partir del lenguaje del propio documento. Si la ambigüedad persiste luego de analizar el texto (\textit{the four corners of the contract}: literalmente, ``las cuatro esquinas'', es decir, todo el documento), el tribunal puede recurrir a las tratativas previas entre las partes (\textit{past dealings}) y a los usos y costumbres del rubro (\textit{relevant industry standards}).

\subsection{Validez versus exigibilidad (\textit{enforceability})}

Un contrato puede ser válido ---contar con todos sus elementos y estar bien celebrado--- y, sin embargo, resultar \textit{unenforceable}: no exigible ante un tribunal.

\begin{example}{Compraventa inmobiliaria no inscripta}
Una compraventa de inmueble no inscripta registralmente es válida entre las partes, pero carece de condiciones formales de oponibilidad frente a terceros. La cuestión de la exigibilidad se plantea recién cuando alguna parte lo lleva ante un tribunal.
\end{example}

% ─────────────────────────────────────────────────────────
\chapter{Figuras precontractuales, análisis de jurisprudencia y estructura de sentencias}
\label{cap:jurisprudencia}
% ─────────────────────────────────────────────────────────

\section{La \textit{invitation to treat}: figura previa a la oferta}
\label{sec:invitation}

En el desarrollo de los elementos de la formación del contrato se presentó la distinción entre oferta (\textit{offer}) y contraoferta (\textit{counter offer}). A estas dos figuras debe agregarse una tercera, que se sitúa en un momento anterior a la oferta misma: la \textbf{\textit{invitation to treat}}.

\begin{definition}{\textit{Invitation to treat}}
\textit{``An invitation to treat is an invitation to start negotiations with intent to create an offer. A solicitation for more offers, usually as a preliminary step to form a contract.''}

La \textit{invitation to treat} es una manera de ``abrir el juego'': decir ``¿qué me proponés?'', a partir de la cual podrá surgir la oferta propiamente dicha. En este momento todavía no existe la oferta como elemento formal del contrato: se trata de una propuesta genérica sin efectos jurídicos vinculantes.
\end{definition}

\begin{example}{La compra de una bicicleta usada}
La publicación en una plataforma de compraventa online constituye una \textit{invitation to treat}: invita a negociar, pero todavía no hay contrato. Cuando el comprador propone un precio, eso es una \textit{offer}. Si el vendedor responde «te lo dejo a tanto», se produce una \textit{counter offer}. Recién cuando una de las partes acepta sin cambios queda formado el contrato.
\end{example}

\begin{example}{La reunión previa a la propuesta}
En contratos comerciales, una reunión destinada a poner en contexto qué se quiere hacer concluye típicamente en el pedido: «hacéme llegar una propuesta». Esa propuesta será la \textit{offer}; la reunión previa es, en cambio, la \textit{invitation to treat}.
\end{example}

La \textit{invitation to treat} guarda semejanza con las tratativas preliminares del derecho argentino. En el derecho argentino, sin embargo, las tratativas no son consideradas un elemento constitutivo del contrato; solo se evalúan si el contrato es objeto de discusión judicial, por ejemplo si existió algún vicio de la voluntad.

\begin{table}[H]
\centering
\caption{Figuras precontractuales y contractuales: \textit{offer}, \textit{counter offer}, \textit{invitation to treat}}
\begin{tabularx}{\textwidth}{>{\raggedright\arraybackslash}p{0.22\textwidth} >{\raggedright\arraybackslash}X >{\raggedright\arraybackslash}X}
\toprule
\textbf{Figura} & \textbf{Qué es} & \textbf{Efecto jurídico} \\
\midrule
\textit{Invitation to treat} & Invitación a iniciar negociaciones con la intención de llegar a una oferta; ``abrir el juego''. & Todavía no existe oferta ni contrato. Etapa precontractual. \\
\textit{Offer} & Manifestación de voluntad de celebrar un contrato bajo términos propuestos, dirigida a quien tiene el poder de aceptarla. & Elemento constitutivo del contrato (junto con la \textit{acceptance}). \\
\textit{Counter offer} & Aceptación condicionada o con cambios en las condiciones. & No es aceptación: la oferta original queda sin efecto y surge una nueva oferta. \\
\textit{Acceptance} & Conformidad con la oferta tal como fue propuesta. & Produce un contrato válido y vinculante (\textit{legally binding}). \\
\bottomrule
\end{tabularx}
\end{table}

\begin{important}
La \textit{invitation to treat} está antes de la oferta: es la etapa de negociaciones previas. La \textit{offer} aceptada tal cual forma el contrato. Si la aceptación introduce cambios, hay una \textit{counter offer} y comienza de nuevo el ciclo.
\end{important}

\section{Análisis del fallo \textit{Hammer v. Sidway} (1891)}
\label{sec:hammer}

\subsection{El caso y los hechos}

El fallo trata de un contrato entre un tío y su sobrino que contiene una obligación de no hacer: el sobrino se comprometió a abstenerse de ciertos vicios (tomar, fumar y otras conductas similares) hasta cumplir la mayoría de edad a los veintiún años. A cambio, el tío prometía una suma de dinero.

El tío y el sobrino comparten el nombre William Story (Senior y Junior, respectivamente). Cuando el sobrino cumplió veintiún años, el tío le confirmó que el dinero estaba depositado en el banco, generando intereses. El tío falleció sin realizar el pago. El sobrino, cansado de esperar, cedió su crédito a un tercero.

\begin{table}[H]
\centering
\caption{Partes y roles en \textit{Hammer v. Sidway}}
\begin{tabularx}{\textwidth}{>{\raggedright\arraybackslash}p{0.22\textwidth} >{\raggedright\arraybackslash}p{0.22\textwidth} >{\raggedright\arraybackslash}X}
\toprule
\textbf{Sujeto} & \textbf{Rol} & \textbf{Explicación} \\
\midrule
Tío (William Story, Senior) & Promitente (\textit{promisor}) & Promete el dinero a cambio de la abstención del sobrino. \\
Sobrino (William Story, Junior) & Beneficiario (\textit{promisee}) & Cumple la abstención; cede su crédito a Hammer. \\
Sidway & Albacea (\textit{executor}) del patrimonio (\textit{estate}) & Ejecuta las disposiciones de última voluntad; se niega a pagar. Es el demandado. \\
Hammer & Tercero acreedor; cesionario (\textit{assignee}) & Reclama el pago al albacea. Es el actor; por eso su nombre encabeza la causa. \\
\bottomrule
\end{tabularx}
\end{table}

\begin{definition}{\textit{Estate}}
El \textit{estate} no es el testamento: es el patrimonio o acervo sucesorio. El testamento (\textit{will}) es el instrumento que distribuye ese patrimonio.
\end{definition}

\begin{note}
No se cede «el contrato» en sí, sino los derechos que surgen de él. En el caso, el sobrino cede a Hammer su derecho a cobrar la promesa del tío. Hammer es cesionario de ese crédito.
\end{note}

\subsection{La cuestión jurídica debatida}

El albacea Sidway se negó a pagar sosteniendo que el contrato era inválido: el tío había prometido dinero, pero el sobrino no había dado nada a cambio, dado que la abstención de una conducta no constituiría una contraprestación válida.

\begin{important}
¿Puede constituir \textit{consideration} (contraprestación válida) el hecho de abstenerse de una conducta que uno tenía derecho a realizar? Esta era la pregunta jurídica central del caso. El albacea sostenía que no: sin contraprestación, no hay contrato.
\end{important}

\subsection{El recorrido procesal y la decisión}

\begin{table}[H]
\centering
\caption{Recorrido procesal del caso}
\begin{tabularx}{\textwidth}{>{\raggedright\arraybackslash}p{0.26\textwidth} >{\raggedright\arraybackslash}p{0.28\textwidth} >{\raggedright\arraybackslash}X}
\toprule
\textbf{Instancia} & \textbf{Resultado} & \textbf{Fundamento} \\
\midrule
Primera instancia & A favor del albacea: no debe pagar. & No hay contraprestación válida; sin ella, no hay contrato. \\
Cámara de Apelaciones de Nueva York & Revoca; falla a favor de Hammer. & La abstención de una conducta lícita es contraprestación válida; no importa cuánto le costó al sobrino renunciar a sus derechos legítimos. \\
\bottomrule
\end{tabularx}
\end{table}

El fundamento de la Cámara fue que el sobrino tenía un derecho legítimo a realizar esas conductas y que dejó de ejercerlo confiando en la promesa de su tío (\textit{upon the faith of his uncle's agreement}). La Cámara agregó que \textit{``we need not speculate''}: no es necesario analizar cuánto le costó al sobrino dejar de fumar o de salir. Basta con advertir que renunció a una libertad legítima.

\begin{definition}{\textit{Forbearance}}
Abstención o autorrestricción: dejar de hacer algo que legítimamente se podría hacer, a cambio de un bien mayor. En el caso, el sobrino restringió su libertad legítima confiando en la promesa de su tío. Es equivalente, en este contexto, a la idea de \textit{refrain from doing something}.
\end{definition}

\subsection{Vocabulario clave del fallo}

\begin{table}[H]
\centering
\caption{Vocabulario central de \textit{Hammer v. Sidway}}
\begin{tabularx}{\textwidth}{>{\raggedright\arraybackslash}p{0.24\textwidth} >{\raggedright\arraybackslash}X}
\toprule
\textbf{Término} & \textbf{Significado en contexto} \\
\midrule
\textit{contend} & Discutir, refutar, debatir. \\
\textit{forbearance} & Abstención; renunciar a algo legítimamente permitido a cambio de un bien mayor. \\
\textit{upon the faith of} & Confiando en la promesa. \\
\textit{indulge} & Permitirse, darse el gusto (en el caso, conductas lícitas como tomar o fumar). \\
\textit{arose} & Pasado de \textit{arise}: surgió. ``The issue before the court arose'': la cuestión ante el tribunal surgió. \\
\textit{assigned} & Cedido. Los derechos que surgen del contrato se ceden, no el contrato mismo. \\
\textit{leading case} & Caso que sienta jurisprudencia y abre el camino interpretativo. \textit{To lead}: guiar, marcar. \\
\bottomrule
\end{tabularx}
\end{table}

\subsection{Vigencia del fallo}

El fallo de 1891 constituye un \textit{leading case} porque, en ese momento, la noción de \textit{consideration} todavía estaba en discusión y se entendía, en general, como el intercambio de cosas materiales o la realización de una conducta positiva. A partir de \textit{Hammer v. Sidway} se interpretó que la contraprestación también puede consistir en un no hacer.

\begin{example}{Los contratos online}
Al aceptar los términos de un servicio online sin pagar, se puede estar renunciando a un derecho ---navegar sin publicidad, sin que se guarden datos de navegación---: esa renuncia funciona como la contraprestación (\textit{consideration}) del acuerdo. Los principios del fallo de 1891 resultan plenamente aplicables a los contratos que se forman actualmente.
\end{example}

\section{Roles procesales en inglés}

\begin{table}[H]
\centering
\caption{Roles procesales en inglés}
\begin{tabularx}{\textwidth}{>{\raggedright\arraybackslash}p{0.20\textwidth} >{\raggedright\arraybackslash}X >{\raggedright\arraybackslash}p{0.24\textwidth}}
\toprule
\textbf{Término} & \textbf{Significado} & \textbf{En el caso analizado} \\
\midrule
\textit{Plaintiff} & Actora, demandante; quien lleva a cabo la acción. & Hammer (tercero acreedor). \\
\textit{Defendant} & Demandado; en materia penal, acusado o imputado. & Sidway (el albacea). \\
\textit{Prosecutor} & El fiscal: cuando quien acciona es el Estado. & No interviene (caso civil). \\
\bottomrule
\end{tabularx}
\end{table}

\section{La estructura de las sentencias en inglés}

Las sentencias en inglés siguen una estructura que, una vez conocida, permite encontrar la información con mayor rapidez que en los fallos argentinos. Los hechos y el derecho aparecen claramente separados y marcados, lo que facilita la investigación jurídica.

\begin{table}[H]
\centering
\caption{Estructura comparada de sentencias: inglés y Argentina}
\begin{tabularx}{\textwidth}{>{\raggedright\arraybackslash}p{0.26\textwidth} >{\raggedright\arraybackslash}X >{\raggedright\arraybackslash}p{0.22\textwidth}}
\toprule
\textbf{Parte en inglés} & \textbf{Contenido} & \textbf{En los fallos argentinos} \\
\midrule
\textit{Case title} / encabezado & Partes, tribunal y jueces; se presenta en un rectángulo. & Encabezado en texto corrido. \\
\textit{Findings of fact} / \textit{facts} & Determinación y narrativa de los hechos pertinentes. & Los hechos y el derecho aparecen algo mezclados. \\
\textit{Determination(s) of law} & Derecho aplicable en sentido amplio: normas, precedentes, principios, usos. & Considerandos. \\
\textit{Application} / \textit{reasoning} & Aplicación del derecho a los hechos; cómo el tribunal conjuga ambos. & Parecido a los votos de cada juez. \\
\textit{Conclusions} / \textit{decision} / \textit{order} & Parte dispositiva o resolutiva. & Parte dispositiva. \\
\bottomrule
\end{tabularx}
\end{table}

\begin{note}
Es el mismo tipo textual (sentencias judiciales), pero estructurado de otra manera. La información se encuentra con mayor facilidad en un fallo en inglés gracias al encabezado en rectángulo y a la separación explícita entre hechos, derecho, análisis y resolución.
\end{note}

\section{Estrategia de lectura de textos jurídicos complejos}

La estrategia propuesta para el inglés jurídico consiste en avanzar de las unidades pequeñas hacia las más grandes: primero extraer la idea de la oración, luego la del párrafo, sin detenerse en cada palabra individual.

La mayoría de las veces es posible evitar el diccionario; corresponde usarlo ante los términos técnicos que condicionan la comprensión global ---como \textit{consideration} o el rol de un tercero como Hammer--- dado que un error en estos términos puede llevar a una comprensión errónea de todo el texto.

Esta propuesta es aplicable no solo al inglés jurídico sino a la lectura de cualquier texto técnico: se trata de desplazar el objetivo de ``entender todo, traducir todo literalmente'' hacia la extracción de ideas generales.

% ─────────────────────────────────────────────────────────
\chapter{Cláusulas estándar en contratos internacionales (\textit{boilerplate clauses})}
\label{cap:boilerplate}
% ─────────────────────────────────────────────────────────

\section{¿Qué son las \textit{boilerplate clauses}?}

\begin{definition}{\textit{Boilerplate clauses}}
Las \textit{boilerplate clauses} (cláusulas estándar o cláusulas tipo) son disposiciones contractuales estandarizadas que se replican con escasas variaciones de un contrato a otro. Debido a su carácter genérico, suelen ser pasadas por alto por las partes al momento de negociar y suscribir un acuerdo. Sin embargo, tienen relevancia jurídica significativa: regulan cuestiones vinculadas a posibles conflictos entre las partes.
\end{definition}

En la práctica argentina, estas cláusulas suelen agruparse bajo el título ``Disposiciones Varias''. Se caracterizan por: ser lo suficientemente estándar para no ser objeto de negociación; abordar cuestiones de conflicto potencial; estar redactadas en lenguaje claro y estandarizado; y ubicarse habitualmente al final del contrato.

\begin{definition}{Origen del término \textit{boilerplate}}
El término proviene de la imprenta del siglo XIX: las placas metálicas (\textit{boilerplates}) contenían textos preestablecidos que se replicaban idénticamente en cada impresión. De allí deriva la idea de texto fijo y no sujeto a modificación. No debe traducirse literalmente; en el lenguaje jurídico significa texto estandarizado, no sujeto a negociación.
\end{definition}

\section{Jurisdicción y derecho aplicable}

\subsection{\textit{Governing law}: el derecho que rige el contrato}

\begin{definition}{\textit{Governing law} / Derecho aplicable}
La cláusula de derecho aplicable determina con qué legislación se interpretará e integrará el contrato. Es especialmente relevante en contratos internacionales, donde las partes provienen de distintas jurisdicciones.
\end{definition}

\subsection{\textit{Jurisdiction}: competencia, no jurisdicción}

La palabra \textit{jurisdiction} constituye un falso cognado parcial: si bien en ocasiones puede referirse efectivamente a la jurisdicción en sentido abstracto, en la mayoría de los contextos contractuales alude a la \textbf{competencia} (qué tribunal específico entiende en la causa).

Las cláusulas de derecho aplicable y de competencia deben ser coherentes entre sí: no resulta válido establecer que el contrato se rige por la legislación de un país y que los conflictos se someterán a los tribunales de otro, si tal combinación carece de coherencia jurídica.

\begin{example}{Inconsistencia entre derecho aplicable y competencia}
Sería incoherente establecer que el contrato se rige por la legislación de Estados Unidos y que cualquier controversia se someterá a los tribunales argentinos. Ambas cláusulas deben articularse de manera consistente.
\end{example}

\section{Riesgo y cumplimiento}

\subsection{\textit{Entire agreement}: el contrato es lo que está escrito}

\begin{definition}{\textit{Entire agreement}}
La cláusula de \textit{entire agreement} establece que únicamente lo consignado por escrito en el contrato es vinculante. Las negociaciones previas, los correos electrónicos y las manifestaciones realizadas durante las tratativas anteriores a la firma \textbf{no} integran el acuerdo ni tienen fuerza obligatoria.
\end{definition}

En el derecho argentino esta cláusula resulta en principio innecesaria, dado que el principio de autonomía de la voluntad es implícito. Sin embargo, por efecto de la globalización comienza a incorporarse en contratos locales adoptados de modelos anglosajones.

\subsection{\textit{Force majeure} y \textit{acts of God}}

\begin{definition}{\textit{Force majeure} / \textit{Acts of God}}
\begin{itemize}
    \item \textbf{\textit{Acts of God}:} Eventos de la naturaleza que ocurren sin intervención humana (terremotos, inundaciones, huracanes).
    \item \textbf{\textit{Force majeure}:} Concepto más amplio que abarca tanto eventos naturales como de origen humano (guerras, disturbios civiles, huelgas, pandemias).
\end{itemize}
Ambos equivalen en términos generales a la fuerza mayor y el caso fortuito del derecho argentino. La pandemia de COVID-19 fue un evento paradigmático de \textit{force majeure} en el ámbito contractual.
\end{definition}

En los contratos redactados en inglés no existe una definición legal predeterminada de fuerza mayor, por lo que la cláusula debe incluirse expresamente y definir con precisión qué circunstancias la configuran. En el derecho argentino, el Código Civil y Comercial ya la define, por lo que no es necesaria una definición exhaustiva en el texto del contrato. El sistema argentino cuenta además con la \textbf{Teoría de la Imprevisión}, que contempla circunstancias sobrevinientes que afectan el cumplimiento normal de las obligaciones.

\subsection{\textit{Terminate} y \textit{discharge}: extinción del contrato}

\begin{definition}{\textit{Terminate} vs. \textit{Discharge}}
\begin{itemize}
    \item \textbf{\textit{Terminate}:} Poner fin al contrato de manera anticipada (extinción anormal). Equivale a las formas anormales de extinción del derecho argentino: rescisión, resolución, revocación, nulidad, novación. \textit{Early termination}: extinción anticipada.
    \item \textbf{\textit{Discharge}:} Concepto más general. Abarca tanto el cumplimiento o vencimiento del plazo (extinción normal) como otras formas de extinción.
\end{itemize}
\end{definition}

\begin{note}
\textit{Termination} puede tener distintos significados según el contexto. En el ámbito laboral, \textit{employee termination} significa despido. La expresión ``\textit{the employee was terminated}'' equivale a ``el empleado fue despedido''. El contexto determina el significado.
\end{note}

\section{Disposiciones operativas}

\subsection{\textit{Assignment} y \textit{novation}: cesión e innovación}

\begin{definition}{\textit{Assignment} / \textit{Novation}}
\begin{itemize}
    \item \textbf{\textit{Assignment} (cesión):} Cesión de los derechos y obligaciones emergentes de un contrato a un tercero. Familia léxica: \textit{assignment} (sustantivo), \textit{assigned} (participio), \textit{to assign} (verbo).
    \item \textbf{\textit{Novation} (novación):} Creación de una nueva obligación que se desprende de una obligación anterior. Mecanismo distinto de la cesión.
\end{itemize}
\end{definition}

\subsection{\textit{Counterparts}: múltiples ejemplares}

\begin{definition}{\textit{Counterparts}}
La cláusula de \textit{counterparts} regula la celebración del contrato mediante múltiples ejemplares firmados posiblemente en lugares distintos. Todos los ejemplares se consideran un único instrumento con idéntico contenido. En el derecho argentino se la incorpora mediante la fórmula clásica: ``dos ejemplares de un mismo tenor y a un solo efecto''.
\end{definition}

\subsection{\textit{Severance clause}: nulidad parcial}

\begin{definition}{\textit{Severance clause}}
Establece que si una cláusula del contrato es declarada nula o inexigible, será removida del contrato sin afectar la validez del resto. El verbo \textit{to sever} significa ``cortar'' o ``separar''. El principio subyacente es el de nulidad parcial.
\end{definition}

\begin{important}
La cláusula tiene un límite: si la cláusula declarada nula es sustancial para la economía del contrato, la nulidad puede afectar al contrato en su totalidad. En el sistema argentino, el Código Civil y Comercial ya contempla el principio de nulidad parcial, por lo que esta cláusula resulta en principio innecesaria.
\end{important}

\section{Disposiciones administrativas}

\subsection{\textit{Relationship of the parties}: naturaleza del vínculo}

Esta cláusula define la naturaleza del vínculo jurídico entre las partes, especificando que no existe relación de dependencia laboral. La mera inserción de esa cláusula, sin embargo, no determina automáticamente la naturaleza del vínculo: si de los hechos surge una relación de dependencia, la calificación jurídica real prevalecerá sobre la designación formal del contrato.

\subsection{\textit{Amendment}: modificación escrita}

La cláusula de \textit{amendment} establece que el contrato solo podrá ser modificado mediante acuerdo escrito de ambas partes: quedan prohibidas las modificaciones unilaterales. En el sistema argentino este principio está implícito en la autonomía de la voluntad.

\subsection{\textit{Notices}: notificaciones formales}

\begin{definition}{\textit{Notices}}
Regula el procedimiento y los medios para las comunicaciones formales con efectos jurídicos: dirección de cada parte, medios admitidos (correo electrónico, fax, servicio de mensajería) y el momento a partir del cual la notificación se tiene por efectuada.
\end{definition}

En el sistema argentino se incorpora en el encabezado del contrato al presentar a las partes, con la indicación del domicilio donde serán válidas todas las notificaciones.

\section{Colocaciones léxicas en el ciclo de vida del contrato}
\label{sec:colocaciones-contratos}

Cada etapa de la vida del contrato tiene asociadas expresiones verbales específicas que no son intercambiables:

\begin{table}[H]
\centering
\caption{Colocaciones léxicas en el ciclo de vida de un contrato}
\begin{tabularx}{\textwidth}{>{\raggedright\arraybackslash}p{0.22\textwidth} >{\raggedright\arraybackslash}p{0.30\textwidth} >{\raggedright\arraybackslash}X}
\toprule
\textbf{Etapa} & \textbf{Expresión en inglés} & \textbf{Equivalente en español} \\
\midrule
Redacción & \textit{Draft a contract} & Redactar un contrato \\
Celebración & \textit{Enter into a contract} & Celebrar / formalizar un contrato \\
 & \textit{Execute an agreement} & Celebrarlo \\
 & \textit{Sign a contract} & Firmar un contrato \\
Cumplimiento & \textit{Perform a contract} & Cumplir / llevar a cabo \\
 & \textit{Comply with} & Cumplir con (una obligación) \\
Exigencia & \textit{Enforce} & Hacer cumplir; exigir el cumplimiento \\
\bottomrule
\end{tabularx}
\end{table}

\begin{important}
El verbo \textit{execute} en inglés jurídico se refiere técnicamente al momento de la \textbf{celebración} del contrato, no al de su ejecución o cumplimiento. Sin embargo, en la práctica aparece utilizado con ambos significados. El contexto determina cuál corresponde.
\end{important}

\section{El impacto de la globalización en la contratación argentina}

La globalización genera una creciente incorporación de cláusulas y estructuras propias del sistema anglosajón en los contratos argentinos: terminología en inglés, adopción de la estructura de cláusulas típicamente anglosajonas, reproducción de modelos sin verificar si el Código Civil y Comercial ya contempla las disposiciones incluidas.

\begin{important}
Los contratos que resultan de mezclar irresponsablemente sistemas jurídicos diferentes han sido denominados informalmente \textbf{``contratos Frankenstein''}: incorporan elementos del derecho anglosajón sin considerar que el sistema argentino ya prevé explícitamente, a través del Código Civil y Comercial, muchas de las cláusulas que en los contratos en inglés deben ser explicitadas. El resultado son redundancias, inconsistencias e incompatibilidades.
\end{important}

La diferencia estructural que explica por qué en el \textit{common law} es necesario explicitar cláusulas que en Argentina son implícitas: en el sistema anglosajón muchos principios no están necesariamente codificados; lo que no consta expresamente en el contrato puede no estar protegido. En el sistema argentino, el Código Civil y Comercial establece los principios generales de los contratos, por lo que muchas de estas cláusulas son implícitas y no es necesario consignarlas.

La recomendación profesional: al redactar o revisar un contrato que incorpora cláusulas de origen anglosajón, es necesario verificar si la disposición ya está contemplada por el ordenamiento argentino, adaptar conscientemente la cláusula al sistema local y mantener la coherencia interna del contrato con el marco normativo aplicable.

% ─────────────────────────────────────────────────────────
% CUADRO CORNELL — PARTE III
% ─────────────────────────────────────────────────────────
\chapter*{Cuadro Cornell: Derecho contractual comparado}
\addcontentsline{toc}{chapter}{Cuadro Cornell: Derecho contractual comparado}
\markboth{Cuadro Cornell: Derecho contractual comparado}{Cuadro Cornell: Derecho contractual comparado}

\begin{longtable}{
    >{\raggedright\arraybackslash}p{0.30\textwidth}
    >{\raggedright\arraybackslash}p{0.60\textwidth}
}
\toprule
\textbf{Preguntas / palabras clave} & \textbf{Notas / respuestas} \\
\midrule
\endfirsthead
\toprule
\textbf{Preguntas / palabras clave} & \textbf{Notas / respuestas} \\
\midrule
\endhead
\textbf{¿Cuáles son los cuatro elementos de la formación del contrato en el \textit{common law}?} &
Intención (\textit{intention}), oferta y aceptación (\textit{offer and acceptance}), contraprestación (\textit{consideration}) y certeza de los términos (\textit{certainty of terms}). Los cuatro son eslabones de una misma cadena. \\

\textbf{¿Cómo se presume la intención en vínculos familiares y comerciales?} &
En vínculos familiares: se presume que NO hay intención de contratar. En vínculos comerciales: se presume que SÍ. Ambas presunciones son refutables con prueba. \\

\textbf{¿Qué resolvió \textit{Merritt v. Merritt} (1970)?} &
Que el vínculo conyugal no impide la existencia de un contrato cuando hay prueba de intención: la señora Merritt probó el acuerdo escrito y el tribunal ordenó la transferencia de la propiedad. \\

\textbf{¿Qué es el \textit{meeting of the minds}?} &
El entendimiento mutuo sobre los términos del contrato (\textit{terms and conditions} = cláusulas). Se logra mediante oferta y aceptación. \\

\textbf{¿Qué es una \textit{counter offer}?} &
Una aceptación condicionada o con cambios: la oferta original queda sin efecto y surge una nueva oferta. Los roles de \textit{offeror} y \textit{offeree} se invierten. \\

\textbf{¿Qué es \textit{consideration} y por qué no se traduce como ``consideración''?} &
Es la contraprestación: el intercambio de promesas o prestaciones recíprocas. ``Consideración'' es un falso amigo. Puede consistir en dinero, bienes, servicios o la abstención de una conducta. \\

\textbf{¿Qué estableció \textit{Hammer v. Sidway} (1891)?} &
Que la abstención de una conducta lícita (\textit{forbearance}) puede constituir contraprestación válida. La Cámara de Nueva York revocó la primera instancia y ordenó el pago al cesionario. Es un \textit{leading case}. \\

\textbf{¿Qué diferencia hay entre \textit{invitation to treat}, \textit{offer} y \textit{counter offer}?} &
\textit{Invitation to treat}: etapa previa a la oferta; invita a negociar, no tiene efectos vinculantes. \textit{Offer}: manifestación de voluntad de contratar bajo términos propuestos. \textit{Counter offer}: aceptación condicionada que deja sin efecto la oferta original. \\

\textbf{¿Cuál es la estructura de una sentencia en inglés?} &
Case title $\to$ findings of fact $\to$ determinations of law $\to$ application/reasoning $\to$ conclusions/decision/order. Los hechos y el derecho aparecen separados, a diferencia de los fallos argentinos. \\

\textbf{¿Qué son las \textit{boilerplate clauses}?} &
Cláusulas contractuales estandarizadas que se replican de contrato en contrato. Regulan cuestiones vinculadas a conflictos; suelen ubicarse al final. Con frecuencia son pasadas por alto durante la negociación. \\

\textbf{¿Qué diferencia hay entre \textit{governing law} y \textit{jurisdiction}?} &
\textit{Governing law}: legislación que rige e interpreta el contrato. \textit{Jurisdiction}: en contexto contractual, la competencia (qué tribunal entiende en los conflictos), no la jurisdicción abstracta. Deben ser coherentes. \\

\textbf{¿Qué diferencia hay entre \textit{terminate} y \textit{discharge}?} &
\textit{Terminate}: extinción anticipada (rescisión, resolución, nulidad). \textit{Discharge}: concepto más general; abarca también la extinción normal por cumplimiento o vencimiento. \\

\textbf{¿Qué es la \textit{severance clause}?} &
La cláusula de nulidad parcial: si una cláusula es declarada nula, se la remueve sin afectar el resto del contrato, salvo que sea sustancial para su economía. \\

\textbf{¿Qué son los ``contratos Frankenstein''?} &
Contratos que incorporan cláusulas anglosajonas en contratos argentinos sin verificar si el Código Civil y Comercial ya las contempla. Generan redundancias, inconsistencias e incompatibilidades. \\

\midrule
\multicolumn{2}{p{0.90\textwidth}}{
\textbf{Resumen:} La Unidad 2 recorrió el derecho contractual del \textit{common law} desde sus bases conceptuales hasta su aplicación práctica en los contratos internacionales actuales. Los cuatro elementos de la formación del contrato ---intención, oferta y aceptación, contraprestación y certeza de los términos--- operan como una cadena en la que cada eslabón es indispensable. La figura de la \textit{invitation to treat} completa el cuadro precontractual. El análisis del fallo \textit{Hammer v. Sidway} demostró que la contraprestación puede consistir en una abstención, principio que hoy se aplica a los contratos online. La estructura de las sentencias en inglés permite ubicar la información con mayor rapidez. Finalmente, las \textit{boilerplate clauses} describen las cláusulas estandarizadas que aparecen en los contratos internacionales, su función en el sistema anglosajón y las diferencias con el sistema argentino, donde muchas de ellas ya están contempladas por el Código Civil y Comercial. La recomendación final: al incorporar cláusulas anglosajonas en contratos argentinos, verificar siempre si el ordenamiento local ya las prevé y adaptar conscientemente al sistema propio.
} \\
\bottomrule
\end{longtable}

\end{document}
